\documentclass[../main.tex]{subfiles}
\begin{document}

\chapter{Introduction of GPU}
\lessoninfo{
  video={M1, lessons 1--5 (Instructor and Course Introduction; Modern Processor Paradigms; How to Write Parallel Programs; Introduction to GPU Architecture; History Trend of GPU Architectures)},
  sources={NVIDIA A100 \cite{nvidia2020ampere} and H100 \cite{nvidia2022hopper} whitepapers (generation figures)},
  textbook={K12 \cite{kim2012gpgpu} ch.~1; A18 \cite{aamodt2018gpu} ch.~1; CPG \cite{nvidia2024cudaguide} §1 Introduction, §2 Programming Model, §4 Hardware Implementation; PMPP \cite{hwu2022pmpp} ch.~1; PPP \cite{mattson2004patterns} ch.~2 (optional); BPG \cite{nvidia2026cudabestpractices} §4.1.3 Understanding Scaling},
  keywords={throughput, latency, virtual ISA, SIMD, SPMD, SIMT, multithreading, Amdahl's law, reduction, Flynn's taxonomy, streaming multiprocessor, warp, unified shader, tensor core},
}

This lesson sets the frame for the whole book.  It contrasts the
design goals of CPUs and GPUs, reviews the techniques by which a CPU
core extracts parallelism (superscalar issue, out-of-order execution,
multithreading, multiprocessing), states Amdahl's law and the two ways of
decomposing a program into parallel work, and then introduces the
vocabulary of GPU hardware: the streaming multiprocessor, the SIMT lane
and the warp.  It closes with the evolution of the GPU from a
fixed-function graphics pipeline to the programmable, tensor-core-equipped
accelerator of today.  Lesson~2 builds on the decomposition ideas with
OpenMP and MPI, and Lesson~3 turns the warp and the SPMD style into
concrete CUDA programs.

\section{CPUs and GPUs}
\label{sec:l01-cpu-vs-gpu}

A \term{GPU}[GPU] (graphics processing unit) was born as a processor
for 3D graphics and has grown into a general parallel processor: today
it is the hardware most associated with machine-learning acceleration and
high-performance computing.\source{M1, CPU vs GPU; K12 §1.2; CPG §1.1 The
Benefits of Using GPUs.}  A \term{CPU}[CPU] (central processing unit) is
built for the opposite point of the design space, and the contrast
between the two, summarised in \cref{tab:l01-cpu-gpu}, organises the
rest of this book.

\begin{table}[htbp]
  \centering
  \caption{CPU and GPU design points.}
  \label{tab:l01-cpu-gpu}
  \begin{tabular}{>{\raggedright\arraybackslash}p{36mm}
                  >{\raggedright\arraybackslash}p{36mm}
                  >{\raggedright\arraybackslash}p{36mm}}
    \toprule
    & CPU & GPU \\
    \midrule
    Target applications & latency-sensitive & throughput-sensitive \\
    Precise exceptions & yes & no \\
    Role & host & accelerator \\
    ISA & public (open) & open or closed; PTX as a public virtual ISA \\
    Programming model & SISD, SIMD & SPMD \\
    \bottomrule
  \end{tabular}
\end{table}

\begin{description}
  \item[Target applications.]  A CPU is optimised for
    \term{latency}[レイテンシ], the time to finish one task, so a
    single thread runs as fast as possible.  A GPU is optimised for
    \term{throughput}[スループット], the amount of work finished per
    unit time, and reaches it by running thousands of threads at once;
    any single thread may be slow.\margin{Latency and throughput are
    independent goals: a design that doubles throughput by running two
    threads side by side leaves each thread's latency unchanged.}
  \item[Precise exceptions.]  A \term{precise exception}[精密例外] is
    an exception after which the architectural state is exactly that of
    the sequential program at the faulting instruction: every earlier
    instruction has completed and no later one has changed state.  CPUs
    guarantee it, because the operating system and debuggers rely on it
    (page faults, breakpoints, signals).  GPUs do not: a faulting kernel
    is normally reported as a whole, and the machine state inside the
    kernel is not recoverable instruction by instruction.
  \item[Host and accelerator.]  The CPU runs the operating system and
    the I/O stack; the GPU does neither, so it depends on the CPU to be
    given work.  The CPU is therefore the \term{host}[ホスト] and the GPU
    the \term{accelerator}[アクセラレータ], a co-processor attached to the
    host over an I/O link.\margin{K12 §1.2 (2012) notes that the PCIe
    link \enquote{provides \SI{4}{\giga\bit\per\second} (Gen~2) or
    \SI{8}{\giga\bit\per\second} (Gen~3)} per lane and
    \enquote{often becomes one of the biggest bottlenecks}
    \cite{kim2012gpgpu}.  Lesson~5 treats host--device transfer.}
  \item[ISA.]  The \term{instruction set architecture (ISA)}[命令セットアーキテクチャ]\index{ISA|see{instruction set architecture (ISA)}}
    of a CPU is public and stable, because compiled software must keep
    running on every later implementation.  A GPU runs in accelerator
    mode behind a driver, which translates the program into whatever
    the current hardware generation executes; the native ISA may be
    closed and may change from one generation to the next.  To give
    programmers a stable target NVIDIA published PTX, a
    \term{virtual ISA}[仮想ISA]: an instruction set that no hardware
    executes directly, but that the driver compiles to the native ISA of
    the installed GPU.\margin{PTX itself, its registers and its place in
    the compiler flow, is the subject of Lesson~10.}
  \item[Programming model.]  CPUs are programmed as
    \term{SISD}[SISD] (single instruction, single data: one instruction
    stream operating on one data item at a time), with
    \term{SIMD}[SIMD] (single instruction, multiple data) vector
    extensions where one instruction applies the same operation to a
    short vector of elements.  GPUs are programmed as
    \term{SPMD}[SPMD] (single program, multiple data): every thread
    runs the same program, but on its own data, and may take its own
    branches.\caution{SPMD describes the \emph{program}, SIMT the
    \emph{hardware} that runs it: a GPU runs SPMD code on SIMT hardware
    (\cref{sec:l01-parallel-programs}); the two acronyms are not
    alternatives.}
\end{description}

SIMD is efficient for regular data-parallel work (image processing,
vector arithmetic) because one instruction fetch drives many operations,
but all lanes execute in lockstep, so data-dependent branches waste
lanes.  SPMD is more flexible, since each thread follows its own control
flow, at the cost of managing many threads and their communication.
Lesson~6 shows how the hardware pays for that flexibility when threads
of one group diverge.

\begin{note}
  The course's own objectives are to explain how a GPU works and to
  optimise programs for it, to prepare for research in GPU
  architecture and compilers, and to explain the recent hardware
  accelerators; it builds a foundation for CUDA rather than a survey of
  its latest features.  It assumes the basic five-stage CPU pipeline and
  C++ programming, and no prior GPU programming.
\end{note}

\section{Modern processor paradigms}
\label{sec:l01-processor-paradigms}

A GPU reuses most of the machinery of a CPU core, so this section
reviews that machinery and the ways in which CPU designers have raised
performance.\source{M1, CPU pipeline and ILP; PMPP ch.~1; A18 ch.~1.}

\subsection{The pipeline}

The CPU model used throughout this book has the stages of
\cref{fig:l01-cpu-pipeline}: a \emph{front end} fetches instructions,
usually from an instruction cache; a \emph{decode} stage interprets them
(a complex x86 instruction is cracked here into several micro-operations);
registers are read after decode; a \emph{scheduler} selects the
instructions to execute; the \emph{execute} stage performs the
computation or the memory access; and a \emph{write-back} stage stores
the result.\margin{Many CPU pipelines also \emph{rename} registers
between decode and scheduling, so that false dependences through reused
register names do not serialise instructions.}  The scheduler is where
the two basic CPU organisations differ.
\begin{description}
  \item[In order.]  A core that uses
    \term{in-order execution}[インオーダ実行] issues instructions in
    program order: an instruction waits until every earlier one has
    issued, even when its own operands are ready.
  \item[Out of order.]  A core that uses
    \term{out-of-order execution}[アウトオブオーダ実行] issues any
    instruction whose operands are ready, regardless of program order,
    so that an instruction stalled on a cache miss does not block the
    independent ones behind it.
\end{description}

\begin{figure}[htbp]
  \centering
  % CPU pipeline timing: single-issue, superscalar and multithreaded (Lesson 1).
% Three charts of a five-stage pipeline (F D S E W).  (a) one instruction
% enters per cycle; (b) a two-wide superscalar enters two per cycle;
% (c) the same two-wide pipeline fed by four threads (one shade per thread).
% Redraws the outline figures m1l2_cpupipeline and m1l2_multi_thread.
% Usage: % CPU pipeline timing: single-issue, superscalar and multithreaded (Lesson 1).
% Three charts of a five-stage pipeline (F D S E W).  (a) one instruction
% enters per cycle; (b) a two-wide superscalar enters two per cycle;
% (c) the same two-wide pipeline fed by four threads (one shade per thread).
% Redraws the outline figures m1l2_cpupipeline and m1l2_multi_thread.
% Usage: % CPU pipeline timing: single-issue, superscalar and multithreaded (Lesson 1).
% Three charts of a five-stage pipeline (F D S E W).  (a) one instruction
% enters per cycle; (b) a two-wide superscalar enters two per cycle;
% (c) the same two-wide pipeline fed by four threads (one shade per thread).
% Redraws the outline figures m1l2_cpupipeline and m1l2_multi_thread.
% Usage: \input{figures/l01-cpu-pipeline}
% Width: body (120 mm); the drawing is about 118 mm wide.
\begin{tikzpicture}[x=3.2mm, y=3.2mm, font=\scriptsize\sffamily]
  % \pipe{x}{y}{fill}: one instruction, five cells from (x,y)
  \newcommand{\pipe}[3]{%
    \foreach \s [count=\i from 0] in {F,D,S,E,W} {
      \draw[fill=#3] (#1+\i,#2) rectangle ++(1,-1)
        node[midway] {\s};
    }}
  % (a) single issue
  \node[anchor=south west, font=\scriptsize\sffamily\bfseries] at (0,0.3)
    {\iflangja{(a) 単一発行}{(a) single issue}};
  \foreach \k in {0,...,3} { \pipe{\k}{-\k}{ln-box} }
  \draw[->, ln-muted] (0,-4.6) -- (8,-4.6)
    node[below left, inner sep=1pt] {\iflangja{サイクル}{cycle}};
  % (b) superscalar, two wide
  \node[anchor=south west, font=\scriptsize\sffamily\bfseries] at (13,0.3)
    {\iflangja{(b) 2命令幅スーパースカラ}{(b) two-wide superscalar}};
  \foreach \k in {0,...,3} {
    \pipe{13+\k}{-2*\k}{ln-box}
    \pipe{13+\k}{-2*\k-1}{ln-box}
  }
  \draw[->, ln-muted] (13,-8.6) -- (21,-8.6)
    node[below left, inner sep=1pt] {\iflangja{サイクル}{cycle}};
  % (c) multithreading on the same two-wide pipeline
  \node[anchor=south west, font=\scriptsize\sffamily\bfseries] at (26,0.3)
    {\iflangja{(c) マルチスレッディング}{(c) multithreading}};
  \foreach \k/\c in {0/ln-accent!25, 1/ln-tip!25, 2/ln-caution!20, 3/ln-source!20} {
    \pipe{26+\k}{-2*\k}{\c}
    \pipe{26+\k}{-2*\k-1}{\c}
  }
  \draw[->, ln-muted] (26,-8.6) -- (34,-8.6)
    node[below left, inner sep=1pt] {\iflangja{サイクル}{cycle}};
  \node[anchor=north west, text=ln-muted, align=left] at (26,-9.3)
    {\iflangja{色＝スレッド}{one shade = one thread}};
\end{tikzpicture}

% Width: body (120 mm); the drawing is about 118 mm wide.
\begin{tikzpicture}[x=3.2mm, y=3.2mm, font=\scriptsize\sffamily]
  % \pipe{x}{y}{fill}: one instruction, five cells from (x,y)
  \newcommand{\pipe}[3]{%
    \foreach \s [count=\i from 0] in {F,D,S,E,W} {
      \draw[fill=#3] (#1+\i,#2) rectangle ++(1,-1)
        node[midway] {\s};
    }}
  % (a) single issue
  \node[anchor=south west, font=\scriptsize\sffamily\bfseries] at (0,0.3)
    {\iflangja{(a) 単一発行}{(a) single issue}};
  \foreach \k in {0,...,3} { \pipe{\k}{-\k}{ln-box} }
  \draw[->, ln-muted] (0,-4.6) -- (8,-4.6)
    node[below left, inner sep=1pt] {\iflangja{サイクル}{cycle}};
  % (b) superscalar, two wide
  \node[anchor=south west, font=\scriptsize\sffamily\bfseries] at (13,0.3)
    {\iflangja{(b) 2命令幅スーパースカラ}{(b) two-wide superscalar}};
  \foreach \k in {0,...,3} {
    \pipe{13+\k}{-2*\k}{ln-box}
    \pipe{13+\k}{-2*\k-1}{ln-box}
  }
  \draw[->, ln-muted] (13,-8.6) -- (21,-8.6)
    node[below left, inner sep=1pt] {\iflangja{サイクル}{cycle}};
  % (c) multithreading on the same two-wide pipeline
  \node[anchor=south west, font=\scriptsize\sffamily\bfseries] at (26,0.3)
    {\iflangja{(c) マルチスレッディング}{(c) multithreading}};
  \foreach \k/\c in {0/ln-accent!25, 1/ln-tip!25, 2/ln-caution!20, 3/ln-source!20} {
    \pipe{26+\k}{-2*\k}{\c}
    \pipe{26+\k}{-2*\k-1}{\c}
  }
  \draw[->, ln-muted] (26,-8.6) -- (34,-8.6)
    node[below left, inner sep=1pt] {\iflangja{サイクル}{cycle}};
  \node[anchor=north west, text=ln-muted, align=left] at (26,-9.3)
    {\iflangja{色＝スレッド}{one shade = one thread}};
\end{tikzpicture}

% Width: body (120 mm); the drawing is about 118 mm wide.
\begin{tikzpicture}[x=3.2mm, y=3.2mm, font=\scriptsize\sffamily]
  % \pipe{x}{y}{fill}: one instruction, five cells from (x,y)
  \newcommand{\pipe}[3]{%
    \foreach \s [count=\i from 0] in {F,D,S,E,W} {
      \draw[fill=#3] (#1+\i,#2) rectangle ++(1,-1)
        node[midway] {\s};
    }}
  % (a) single issue
  \node[anchor=south west, font=\scriptsize\sffamily\bfseries] at (0,0.3)
    {\iflangja{(a) 単一発行}{(a) single issue}};
  \foreach \k in {0,...,3} { \pipe{\k}{-\k}{ln-box} }
  \draw[->, ln-muted] (0,-4.6) -- (8,-4.6)
    node[below left, inner sep=1pt] {\iflangja{サイクル}{cycle}};
  % (b) superscalar, two wide
  \node[anchor=south west, font=\scriptsize\sffamily\bfseries] at (13,0.3)
    {\iflangja{(b) 2命令幅スーパースカラ}{(b) two-wide superscalar}};
  \foreach \k in {0,...,3} {
    \pipe{13+\k}{-2*\k}{ln-box}
    \pipe{13+\k}{-2*\k-1}{ln-box}
  }
  \draw[->, ln-muted] (13,-8.6) -- (21,-8.6)
    node[below left, inner sep=1pt] {\iflangja{サイクル}{cycle}};
  % (c) multithreading on the same two-wide pipeline
  \node[anchor=south west, font=\scriptsize\sffamily\bfseries] at (26,0.3)
    {\iflangja{(c) マルチスレッディング}{(c) multithreading}};
  \foreach \k/\c in {0/ln-accent!25, 1/ln-tip!25, 2/ln-caution!20, 3/ln-source!20} {
    \pipe{26+\k}{-2*\k}{\c}
    \pipe{26+\k}{-2*\k-1}{\c}
  }
  \draw[->, ln-muted] (26,-8.6) -- (34,-8.6)
    node[below left, inner sep=1pt] {\iflangja{サイクル}{cycle}};
  \node[anchor=north west, text=ln-muted, align=left] at (26,-9.3)
    {\iflangja{色＝スレッド}{one shade = one thread}};
\end{tikzpicture}

  \caption{A five-stage pipeline (fetch, decode, schedule, execute,
    write-back) fed with one instruction per cycle (a), with two
    instructions per cycle in a two-wide superscalar (b), and with two
    instructions per cycle drawn from four threads (c).}
  \label{fig:l01-cpu-pipeline}
\end{figure}

\subsection{Instruction-level parallelism}

A single-issue pipeline starts at most one instruction per cycle.  A
\term{superscalar}[スーパースカラ] processor fetches, decodes and
executes several instructions per cycle (\cref{fig:l01-cpu-pipeline}b),
so its \term{instructions per cycle (IPC)}[サイクル当たり命令数]\index{IPC|see{instructions per cycle (IPC)}}
exceeds one.\margin{Width in practice: Intel's Golden Cove core (2021)
decodes up to six x86 instructions per cycle
\cite[§2.4.1]{intel2026optimization}.}  How many instructions can actually run
together is bounded by the program's
\term{instruction-level parallelism (ILP)}[命令レベル並列性]\index{ILP|see{instruction-level parallelism (ILP)}}: the
number of instructions that are independent of each other and could
execute at the same time.

\begin{example}
  \label{ex:l01-ilp}
  A machine has two execution units.  In the pair
  \begin{center}
    \begin{tabular}{ll}
      \toprule
      (a) independent & (b) dependent \\
      \midrule
      \texttt{R1 = R2 * R3} & \texttt{R1 = R2 * R3} \\
      \texttt{R4 = R5 - R6} & \texttt{R4 = R1 - R6} \\
      \bottomrule
    \end{tabular}
  \end{center}
  the instructions of (a) read and write disjoint registers, so both
  issue in the same cycle.  In (b) the second instruction reads
  \texttt{R1}, which the first one writes, so it cannot start before the
  first has produced its result: the pair takes two issue slots in
  sequence and the second unit stays idle.\tip{Two instructions are
  independent when neither reads or writes a register the other
  writes.  Out-of-order hardware looks further down the stream for an
  instruction that passes this test.}
\end{example}

Out-of-order execution exists to find more independent instructions
than the next one in program order, and raising ILP has been one of the
main themes of CPU design.  Two further levers raise single-thread
performance without more parallelism: a deeper pipeline allows a higher
clock frequency, but needs better \term{branch prediction}[分岐予測],
guessing the direction and target of a branch before it executes so
that the longer pipeline stays full; and larger caches cut the number of
memory accesses that miss and reach DRAM.\margin{The Pentium~4 took the
first lever to 20 pipeline stages (2000) \cite{hinton2001pentium4} and
31 (2004) \cite{boggs2004pentium4}.}

\subsection{Multithreading and multiprocessors}

The next source of parallelism is \term{multithreading}[マルチスレッディング]:
a core holds the state of several threads at once and, when the current
thread has no ready instruction, issues from another thread
(\cref{fig:l01-cpu-pipeline}c).  Independent threads supply ready
instructions without any search for ILP, and the hardware cost is
modest, chiefly the extra program counters and registers.
Multithreading is one of the key ingredients of GPU performance: a GPU
core keeps tens of thread groups in flight and switches among them every
cycle.\margin{On the GPU this is fine-grained multithreading with a
context switch every cycle; Lesson~4 shows how it hides memory latency.}

The last lever is to replicate whole cores: a
\term{multiprocessor}[マルチプロセッサ] contains several cores, each
with its own pipeline and usually its own first-level cache, and raises
throughput in proportion to the number of cores for work that can be
divided among them.  Both CPUs and GPUs are multiprocessors; they differ
in how they spend their transistors (\cref{fig:l01-cpu-vs-gpu}).  A CPU
has a few large cores, each built for ILP and single-thread latency,
over a large shared cache.  A GPU has many small cores, each
multithreaded and SIMD-wide, with a comparatively small cache between
them.\margin{Scale today: AMD's EPYC~9654 CPU has 96 cores (2022)
\cite{amd2026epyc9654}; NVIDIA's H100 has 132 streaming multiprocessors
\cite{nvidia2022hopper}, each with 128 lanes.}

\begin{figure}[htbp]
  \centering
  % CPU vs GPU chip organisation (Lesson 1).
% Left: a CPU with four large cores over one large cache.  Right: a GPU
% with sixteen small cores and a thin cache between the rows.
% Redraws the outline figure m1l2_cpu_gpu (also used for m1l4_gpuarch).
% Usage: % CPU vs GPU chip organisation (Lesson 1).
% Left: a CPU with four large cores over one large cache.  Right: a GPU
% with sixteen small cores and a thin cache between the rows.
% Redraws the outline figure m1l2_cpu_gpu (also used for m1l4_gpuarch).
% Usage: % CPU vs GPU chip organisation (Lesson 1).
% Left: a CPU with four large cores over one large cache.  Right: a GPU
% with sixteen small cores and a thin cache between the rows.
% Redraws the outline figure m1l2_cpu_gpu (also used for m1l4_gpuarch).
% Usage: \input{figures/l01-cpu-vs-gpu}
% Width: body (120 mm); the drawing is about 100 mm wide.
\begin{tikzpicture}[x=1mm, y=1mm, font=\scriptsize\sffamily]
  % CPU
  \node[font=\small\sffamily\bfseries] at (20,34) {CPU};
  \foreach \i in {0,1} \foreach \j in {0,1} {
    \draw[fill=ln-box] (\i*20,\j*13+5) rectangle ++(19,12)
      node[midway] {\iflangja{コア}{Core}};
  }
  \draw[fill=ln-accent!15] (0,0) rectangle (39,4)
    node[midway] {\iflangja{キャッシュ}{Cache}};
  % GPU
  \node[font=\small\sffamily\bfseries] at (80,34) {GPU};
  \foreach \i in {0,...,3} \foreach \j in {0,1} {
    \draw[fill=ln-box] (60+\i*10,\j*6+19) rectangle ++(9,5)
      node[midway] {\iflangja{コア}{Core}};
  }
  \draw[fill=ln-accent!15] (60,15) rectangle (99,18)
    node[midway] {\iflangja{キャッシュ}{Cache}};
  \foreach \i in {0,...,3} \foreach \j in {0,1} {
    \draw[fill=ln-box] (60+\i*10,\j*6+3) rectangle ++(9,5)
      node[midway] {\iflangja{コア}{Core}};
  }
  \draw[fill=ln-accent!15] (60,0) rectangle (99,2) ;
\end{tikzpicture}

% Width: body (120 mm); the drawing is about 100 mm wide.
\begin{tikzpicture}[x=1mm, y=1mm, font=\scriptsize\sffamily]
  % CPU
  \node[font=\small\sffamily\bfseries] at (20,34) {CPU};
  \foreach \i in {0,1} \foreach \j in {0,1} {
    \draw[fill=ln-box] (\i*20,\j*13+5) rectangle ++(19,12)
      node[midway] {\iflangja{コア}{Core}};
  }
  \draw[fill=ln-accent!15] (0,0) rectangle (39,4)
    node[midway] {\iflangja{キャッシュ}{Cache}};
  % GPU
  \node[font=\small\sffamily\bfseries] at (80,34) {GPU};
  \foreach \i in {0,...,3} \foreach \j in {0,1} {
    \draw[fill=ln-box] (60+\i*10,\j*6+19) rectangle ++(9,5)
      node[midway] {\iflangja{コア}{Core}};
  }
  \draw[fill=ln-accent!15] (60,15) rectangle (99,18)
    node[midway] {\iflangja{キャッシュ}{Cache}};
  \foreach \i in {0,...,3} \foreach \j in {0,1} {
    \draw[fill=ln-box] (60+\i*10,\j*6+3) rectangle ++(9,5)
      node[midway] {\iflangja{コア}{Core}};
  }
  \draw[fill=ln-accent!15] (60,0) rectangle (99,2) ;
\end{tikzpicture}

% Width: body (120 mm); the drawing is about 100 mm wide.
\begin{tikzpicture}[x=1mm, y=1mm, font=\scriptsize\sffamily]
  % CPU
  \node[font=\small\sffamily\bfseries] at (20,34) {CPU};
  \foreach \i in {0,1} \foreach \j in {0,1} {
    \draw[fill=ln-box] (\i*20,\j*13+5) rectangle ++(19,12)
      node[midway] {\iflangja{コア}{Core}};
  }
  \draw[fill=ln-accent!15] (0,0) rectangle (39,4)
    node[midway] {\iflangja{キャッシュ}{Cache}};
  % GPU
  \node[font=\small\sffamily\bfseries] at (80,34) {GPU};
  \foreach \i in {0,...,3} \foreach \j in {0,1} {
    \draw[fill=ln-box] (60+\i*10,\j*6+19) rectangle ++(9,5)
      node[midway] {\iflangja{コア}{Core}};
  }
  \draw[fill=ln-accent!15] (60,15) rectangle (99,18)
    node[midway] {\iflangja{キャッシュ}{Cache}};
  \foreach \i in {0,...,3} \foreach \j in {0,1} {
    \draw[fill=ln-box] (60+\i*10,\j*6+3) rectangle ++(9,5)
      node[midway] {\iflangja{コア}{Core}};
  }
  \draw[fill=ln-accent!15] (60,0) rectangle (99,2) ;
\end{tikzpicture}

  \caption{A CPU spends its area on a few large cores and a large cache;
    a GPU on many small cores and a small cache.}
  \label{fig:l01-cpu-vs-gpu}
\end{figure}

\begin{keypoint}
  CPU performance comes from ILP (superscalar, out-of-order), clock
  frequency (deep pipelines plus branch prediction), caches,
  multithreading and multiple cores.  A GPU keeps multithreading and
  multiple cores, drops most of the ILP machinery, and multiplies the
  thread count instead.
\end{keypoint}

\section{How to write parallel programs}
\label{sec:l01-parallel-programs}

Hardware with many cores only pays off for programs that expose parallel
work.  This section gives the law that bounds the gain, the two ways of
splitting a program into parallel parts, and the taxonomy in which GPUs
are placed.\source{M1, Amdahl's law and decomposition; PPP ch.~2
\cite{mattson2004patterns}; BPG §4.1.3 Understanding Scaling
\cite{nvidia2026cudabestpractices}.}

\subsection{Amdahl's law}

\begin{definition}[Amdahl's law]
  \label{def:l01-amdahl}
  Let a program spend a fraction $S$ of its sequential execution time in
  work that cannot be parallelised, the \term{serial fraction}[逐次部分の割合],
  and the remaining fraction $P = 1 - S$ in work that $N$ processors
  share evenly.  \term{Amdahl's law}[アムダールの法則] gives the
  speedup over the sequential execution as\margin{Notation: $S$ and
  $P$ are fractions of the \emph{sequential} run time, $S + P = 1$;
  $N$ is the number of processors; the parallel part is assumed to
  divide perfectly.}
  \begin{equation}
    \text{speedup} \;=\; \frac{1}{\,S + P/N\,}.
    \label{eq:l01-amdahl}
  \end{equation}
\end{definition}

When the serial fraction is small and the parallel part is large, the
speedup grows almost linearly with $N$.  When the serial fraction is
significant, scalability collapses: as $N \to \infty$ the term $P/N$
vanishes and the speedup approaches $1/S$, however much hardware is
added.\tip{The ceiling is $1/S$.  A serial fraction of \SI{10}{\percent}
caps the speedup at $10\times$; \SI{1}{\percent} at $100\times$.}

\begin{example}
  \label{ex:l01-amdahl}
  A program has a serial fraction $S = 0.1$.  With $N = 8$ processors
  \cref{eq:l01-amdahl} gives $1/(0.1 + 0.9/8) = 1/0.2125 \approx 4.7$;
  with $N = 64$ it gives $1/(0.1 + 0.9/64) \approx 8.8$; with
  $N = 1024$, $1/(0.1 + 0.9/1024) \approx 9.9$.  Going from 8 to 1024
  processors, a $128\times$ increase in hardware, barely doubles the
  speedup, because the serial \SI{10}{\percent} alone already costs as
  much as the whole parallel part on 8 processors.
\end{example}

\Cref{fig:l01-amdahl} plots \cref{eq:l01-amdahl} against the serial
fraction for several $N$.  The curves stay bunched together until $S$
falls below about $0.3$ and only separate for small serial fractions,
which is why the first task of a parallel programmer is to shrink the
serial part, not to add processors.\caution{Amdahl's law assumes a
fixed problem size (strong scaling).  Gustafson's law, BPG §4.1.3.2
Weak Scaling and Gustafson's Law \cite{nvidia2026cudabestpractices},
grows the problem with $N$ and gives $N + (1 - N)\,s$, where $s$ is
the serial fraction of the \emph{parallel} run: the serial part then
matters less.  The BPG names meshes, grids and some Monte Carlo
simulations as workloads that grow with $N$.}

\begin{figure}[htbp]
  \centering
  % Amdahl's law: speedup against the serial fraction S for N = 2, 4, 8, 16
% and 100 processors (Lesson 1).  speedup = 1 / (S + (1 - S)/N), plotted
% for S from 0.05 to 1.  Redraws the outline figure m1l3_amdahllaw with the
% serial fraction increasing to the right.
% Usage: % Amdahl's law: speedup against the serial fraction S for N = 2, 4, 8, 16
% and 100 processors (Lesson 1).  speedup = 1 / (S + (1 - S)/N), plotted
% for S from 0.05 to 1.  Redraws the outline figure m1l3_amdahllaw with the
% serial fraction increasing to the right.
% Usage: % Amdahl's law: speedup against the serial fraction S for N = 2, 4, 8, 16
% and 100 processors (Lesson 1).  speedup = 1 / (S + (1 - S)/N), plotted
% for S from 0.05 to 1.  Redraws the outline figure m1l3_amdahllaw with the
% serial fraction increasing to the right.
% Usage: \input{figures/l01-amdahl}
% Width: body (120 mm); the drawing is about 95 mm wide.
\begin{tikzpicture}[x=80mm, y=2.6mm, font=\scriptsize\sffamily,
                    >={Stealth[length=1.4mm]}]
  % grid and axes
  \foreach \yv in {2,4,...,16} {
    \draw[ln-rule!50, thin] (0,\yv) -- (1,\yv);
    \node[left, inner sep=1.5pt] at (0,\yv) {\yv};
  }
  \draw[->] (0,0) -- (1.04,0);
  \node[anchor=north, inner sep=1pt] at (0.52,-2)
    {\iflangja{逐次部分の割合 $S$}{serial fraction $S$}};
  \draw[->] (0,0) -- (0,18.5)
    node[above, inner sep=2pt] {\iflangja{高速化率}{speedup}};
  \foreach \xv/\xl in {0.1/0.1, 0.2/0.2, 0.3/0.3, 0.4/0.4, 0.5/0.5,
                       0.6/0.6, 0.7/0.7, 0.8/0.8, 0.9/0.9, 1.0/1.0} {
    \draw (\xv,0) -- (\xv,-0.4) node[below, inner sep=1pt] {\xl};
  }
  % curves
  \draw[thick, ln-accent] plot coordinates {
    (0.05,16.81) (0.10,9.17) (0.15,6.31) (0.20,4.81) (0.25,3.88) (0.30,3.26)
    (0.35,2.81) (0.40,2.46) (0.45,2.20) (0.50,1.98) (0.55,1.80) (0.60,1.66)
    (0.65,1.53) (0.70,1.42) (0.75,1.33) (0.80,1.25) (0.85,1.17) (0.90,1.11)
    (0.95,1.05) (1.00,1.00)};
  \node[right, text=ln-accent, inner sep=1pt] at (0.05,16.81) {$N=100$};
  \draw[thick, ln-source] plot coordinates {
    (0.05,9.14) (0.10,6.40) (0.15,4.92) (0.20,4.00) (0.25,3.37) (0.30,2.91)
    (0.35,2.56) (0.40,2.29) (0.45,2.06) (0.50,1.88) (0.55,1.73) (0.60,1.60)
    (0.65,1.49) (0.70,1.39) (0.75,1.31) (0.80,1.23) (0.85,1.16) (0.90,1.10)
    (0.95,1.05) (1.00,1.00)};
  \node[right, text=ln-source, inner sep=1pt] at (0.05,9.14) {$N=16$};
  \draw[thick, ln-caution] plot coordinates {
    (0.05,5.93) (0.10,4.71) (0.15,3.90) (0.20,3.33) (0.25,2.91) (0.30,2.58)
    (0.35,2.32) (0.40,2.11) (0.45,1.93) (0.50,1.78) (0.55,1.65) (0.60,1.54)
    (0.65,1.44) (0.70,1.36) (0.75,1.28) (0.80,1.21) (0.85,1.15) (0.90,1.10)
    (0.95,1.05) (1.00,1.00)};
  \node[right, text=ln-caution, inner sep=1pt] at (0.05,5.93) {$N=8$};
  \draw[thick, ln-tip] plot coordinates {
    (0.05,3.48) (0.10,3.08) (0.15,2.76) (0.20,2.50) (0.25,2.29) (0.30,2.11)
    (0.35,1.95) (0.40,1.82) (0.45,1.70) (0.50,1.60) (0.55,1.51) (0.60,1.43)
    (0.65,1.36) (0.70,1.29) (0.75,1.23) (0.80,1.18) (0.85,1.13) (0.90,1.08)
    (0.95,1.04) (1.00,1.00)};
  \node[right, text=ln-tip, inner sep=1pt] at (0.05,3.48) {$N=4$};
  \draw[thick, ln-muted] plot coordinates {
    (0.05,1.90) (0.10,1.82) (0.15,1.74) (0.20,1.67) (0.25,1.60) (0.30,1.54)
    (0.35,1.48) (0.40,1.43) (0.45,1.38) (0.50,1.33) (0.55,1.29) (0.60,1.25)
    (0.65,1.21) (0.70,1.18) (0.75,1.14) (0.80,1.11) (0.85,1.08) (0.90,1.05)
    (0.95,1.03) (1.00,1.00)};
  \node[right, text=ln-muted, inner sep=1pt] at (0.05,1.90) {$N=2$};
  % 1/S bound
  \draw[densely dashed, ln-muted] (0.1,10)
    node[above right, inner sep=1pt, text=ln-muted]
    {\iflangja{$S=0.1$: 上限 $1/S=10$}{$S=0.1$: bound $1/S=10$}} -- (0.1,0);
\end{tikzpicture}

% Width: body (120 mm); the drawing is about 95 mm wide.
\begin{tikzpicture}[x=80mm, y=2.6mm, font=\scriptsize\sffamily,
                    >={Stealth[length=1.4mm]}]
  % grid and axes
  \foreach \yv in {2,4,...,16} {
    \draw[ln-rule!50, thin] (0,\yv) -- (1,\yv);
    \node[left, inner sep=1.5pt] at (0,\yv) {\yv};
  }
  \draw[->] (0,0) -- (1.04,0);
  \node[anchor=north, inner sep=1pt] at (0.52,-2)
    {\iflangja{逐次部分の割合 $S$}{serial fraction $S$}};
  \draw[->] (0,0) -- (0,18.5)
    node[above, inner sep=2pt] {\iflangja{高速化率}{speedup}};
  \foreach \xv/\xl in {0.1/0.1, 0.2/0.2, 0.3/0.3, 0.4/0.4, 0.5/0.5,
                       0.6/0.6, 0.7/0.7, 0.8/0.8, 0.9/0.9, 1.0/1.0} {
    \draw (\xv,0) -- (\xv,-0.4) node[below, inner sep=1pt] {\xl};
  }
  % curves
  \draw[thick, ln-accent] plot coordinates {
    (0.05,16.81) (0.10,9.17) (0.15,6.31) (0.20,4.81) (0.25,3.88) (0.30,3.26)
    (0.35,2.81) (0.40,2.46) (0.45,2.20) (0.50,1.98) (0.55,1.80) (0.60,1.66)
    (0.65,1.53) (0.70,1.42) (0.75,1.33) (0.80,1.25) (0.85,1.17) (0.90,1.11)
    (0.95,1.05) (1.00,1.00)};
  \node[right, text=ln-accent, inner sep=1pt] at (0.05,16.81) {$N=100$};
  \draw[thick, ln-source] plot coordinates {
    (0.05,9.14) (0.10,6.40) (0.15,4.92) (0.20,4.00) (0.25,3.37) (0.30,2.91)
    (0.35,2.56) (0.40,2.29) (0.45,2.06) (0.50,1.88) (0.55,1.73) (0.60,1.60)
    (0.65,1.49) (0.70,1.39) (0.75,1.31) (0.80,1.23) (0.85,1.16) (0.90,1.10)
    (0.95,1.05) (1.00,1.00)};
  \node[right, text=ln-source, inner sep=1pt] at (0.05,9.14) {$N=16$};
  \draw[thick, ln-caution] plot coordinates {
    (0.05,5.93) (0.10,4.71) (0.15,3.90) (0.20,3.33) (0.25,2.91) (0.30,2.58)
    (0.35,2.32) (0.40,2.11) (0.45,1.93) (0.50,1.78) (0.55,1.65) (0.60,1.54)
    (0.65,1.44) (0.70,1.36) (0.75,1.28) (0.80,1.21) (0.85,1.15) (0.90,1.10)
    (0.95,1.05) (1.00,1.00)};
  \node[right, text=ln-caution, inner sep=1pt] at (0.05,5.93) {$N=8$};
  \draw[thick, ln-tip] plot coordinates {
    (0.05,3.48) (0.10,3.08) (0.15,2.76) (0.20,2.50) (0.25,2.29) (0.30,2.11)
    (0.35,1.95) (0.40,1.82) (0.45,1.70) (0.50,1.60) (0.55,1.51) (0.60,1.43)
    (0.65,1.36) (0.70,1.29) (0.75,1.23) (0.80,1.18) (0.85,1.13) (0.90,1.08)
    (0.95,1.04) (1.00,1.00)};
  \node[right, text=ln-tip, inner sep=1pt] at (0.05,3.48) {$N=4$};
  \draw[thick, ln-muted] plot coordinates {
    (0.05,1.90) (0.10,1.82) (0.15,1.74) (0.20,1.67) (0.25,1.60) (0.30,1.54)
    (0.35,1.48) (0.40,1.43) (0.45,1.38) (0.50,1.33) (0.55,1.29) (0.60,1.25)
    (0.65,1.21) (0.70,1.18) (0.75,1.14) (0.80,1.11) (0.85,1.08) (0.90,1.05)
    (0.95,1.03) (1.00,1.00)};
  \node[right, text=ln-muted, inner sep=1pt] at (0.05,1.90) {$N=2$};
  % 1/S bound
  \draw[densely dashed, ln-muted] (0.1,10)
    node[above right, inner sep=1pt, text=ln-muted]
    {\iflangja{$S=0.1$: 上限 $1/S=10$}{$S=0.1$: bound $1/S=10$}} -- (0.1,0);
\end{tikzpicture}

% Width: body (120 mm); the drawing is about 95 mm wide.
\begin{tikzpicture}[x=80mm, y=2.6mm, font=\scriptsize\sffamily,
                    >={Stealth[length=1.4mm]}]
  % grid and axes
  \foreach \yv in {2,4,...,16} {
    \draw[ln-rule!50, thin] (0,\yv) -- (1,\yv);
    \node[left, inner sep=1.5pt] at (0,\yv) {\yv};
  }
  \draw[->] (0,0) -- (1.04,0);
  \node[anchor=north, inner sep=1pt] at (0.52,-2)
    {\iflangja{逐次部分の割合 $S$}{serial fraction $S$}};
  \draw[->] (0,0) -- (0,18.5)
    node[above, inner sep=2pt] {\iflangja{高速化率}{speedup}};
  \foreach \xv/\xl in {0.1/0.1, 0.2/0.2, 0.3/0.3, 0.4/0.4, 0.5/0.5,
                       0.6/0.6, 0.7/0.7, 0.8/0.8, 0.9/0.9, 1.0/1.0} {
    \draw (\xv,0) -- (\xv,-0.4) node[below, inner sep=1pt] {\xl};
  }
  % curves
  \draw[thick, ln-accent] plot coordinates {
    (0.05,16.81) (0.10,9.17) (0.15,6.31) (0.20,4.81) (0.25,3.88) (0.30,3.26)
    (0.35,2.81) (0.40,2.46) (0.45,2.20) (0.50,1.98) (0.55,1.80) (0.60,1.66)
    (0.65,1.53) (0.70,1.42) (0.75,1.33) (0.80,1.25) (0.85,1.17) (0.90,1.11)
    (0.95,1.05) (1.00,1.00)};
  \node[right, text=ln-accent, inner sep=1pt] at (0.05,16.81) {$N=100$};
  \draw[thick, ln-source] plot coordinates {
    (0.05,9.14) (0.10,6.40) (0.15,4.92) (0.20,4.00) (0.25,3.37) (0.30,2.91)
    (0.35,2.56) (0.40,2.29) (0.45,2.06) (0.50,1.88) (0.55,1.73) (0.60,1.60)
    (0.65,1.49) (0.70,1.39) (0.75,1.31) (0.80,1.23) (0.85,1.16) (0.90,1.10)
    (0.95,1.05) (1.00,1.00)};
  \node[right, text=ln-source, inner sep=1pt] at (0.05,9.14) {$N=16$};
  \draw[thick, ln-caution] plot coordinates {
    (0.05,5.93) (0.10,4.71) (0.15,3.90) (0.20,3.33) (0.25,2.91) (0.30,2.58)
    (0.35,2.32) (0.40,2.11) (0.45,1.93) (0.50,1.78) (0.55,1.65) (0.60,1.54)
    (0.65,1.44) (0.70,1.36) (0.75,1.28) (0.80,1.21) (0.85,1.15) (0.90,1.10)
    (0.95,1.05) (1.00,1.00)};
  \node[right, text=ln-caution, inner sep=1pt] at (0.05,5.93) {$N=8$};
  \draw[thick, ln-tip] plot coordinates {
    (0.05,3.48) (0.10,3.08) (0.15,2.76) (0.20,2.50) (0.25,2.29) (0.30,2.11)
    (0.35,1.95) (0.40,1.82) (0.45,1.70) (0.50,1.60) (0.55,1.51) (0.60,1.43)
    (0.65,1.36) (0.70,1.29) (0.75,1.23) (0.80,1.18) (0.85,1.13) (0.90,1.08)
    (0.95,1.04) (1.00,1.00)};
  \node[right, text=ln-tip, inner sep=1pt] at (0.05,3.48) {$N=4$};
  \draw[thick, ln-muted] plot coordinates {
    (0.05,1.90) (0.10,1.82) (0.15,1.74) (0.20,1.67) (0.25,1.60) (0.30,1.54)
    (0.35,1.48) (0.40,1.43) (0.45,1.38) (0.50,1.33) (0.55,1.29) (0.60,1.25)
    (0.65,1.21) (0.70,1.18) (0.75,1.14) (0.80,1.11) (0.85,1.08) (0.90,1.05)
    (0.95,1.03) (1.00,1.00)};
  \node[right, text=ln-muted, inner sep=1pt] at (0.05,1.90) {$N=2$};
  % 1/S bound
  \draw[densely dashed, ln-muted] (0.1,10)
    node[above right, inner sep=1pt, text=ln-muted]
    {\iflangja{$S=0.1$: 上限 $1/S=10$}{$S=0.1$: bound $1/S=10$}} -- (0.1,0);
\end{tikzpicture}

  \caption{Speedup under Amdahl's law as a function of the serial
    fraction $S$ for $N = 2$, $4$, $8$, $16$ and $100$ processors.  At
    $S = 0.1$ even $N = 100$ gives only $9.2\times$.}
  \label{fig:l01-amdahl}
\end{figure}

\subsection{Task and data decomposition}

Given a program with enough parallel work, there are two ways of
distributing it.  The running example is an array $A$ over which both
the sum and the minimum and maximum are wanted, on two cores
(\cref{fig:l01-decomposition}).
\begin{description}
  \item[By task.]  The split called
    \term{task decomposition}[タスク分割] gives different
    \emph{operations} to different cores: core~1 computes the sum, core~2
    the minimum and maximum, and the whole of $A$ must be sent to both.
  \item[By data.]  The split called
    \term{data decomposition}[データ分割] gives different \emph{data}
    to the cores, each of which performs every operation on
    its part: core~1 sums and scans the first half of $A$, core~2 the
    second half.  The partial results then have to be combined; the
    combining step is a \term{reduction}[リダクション], an operation
    that folds many values into one with an associative operator (sum,
    minimum, maximum).\margin{Lesson~2 shows the reduction as an OpenMP
    clause and as an MPI collective; Lesson~5 touches the
    warp-level primitives that implement it on a GPU.}
\end{description}
Data decomposition scales naturally: with more cores each simply takes
a smaller part of $A$, while task decomposition is limited by the number
of distinct operations.  Its price is the reduction, which is serial
work in the sense of Amdahl's law unless it is itself parallelised as a
tree.

\begin{figure}[htbp]
  \centering
  % Task decomposition versus data decomposition (Lesson 1).
% (a) Task decomposition: the whole array goes to both cores, one sums and
% the other finds min/max.  (b) Data decomposition: each core gets half the
% array, does both operations, and a reduction combines the partial results.
% Redraws the outline figures m1l3_taskdecompose and m1l3_datadecompose.
% Usage: % Task decomposition versus data decomposition (Lesson 1).
% (a) Task decomposition: the whole array goes to both cores, one sums and
% the other finds min/max.  (b) Data decomposition: each core gets half the
% array, does both operations, and a reduction combines the partial results.
% Redraws the outline figures m1l3_taskdecompose and m1l3_datadecompose.
% Usage: % Task decomposition versus data decomposition (Lesson 1).
% (a) Task decomposition: the whole array goes to both cores, one sums and
% the other finds min/max.  (b) Data decomposition: each core gets half the
% array, does both operations, and a reduction combines the partial results.
% Redraws the outline figures m1l3_taskdecompose and m1l3_datadecompose.
% Usage: \input{figures/l01-decomposition}
% Width: body (120 mm); the drawing is about 118 mm wide.
\begin{tikzpicture}[x=1mm, y=1mm, font=\scriptsize\sffamily,
    >={Stealth[length=1.6mm]},
    core/.style={draw, rounded corners=1pt, fill=ln-box, minimum width=12mm,
                 minimum height=5mm},
    res/.style={draw, fill=ln-accent!12, minimum width=16mm, minimum height=5mm,
                align=center}]
  % (a) task decomposition
  \node[anchor=west, font=\scriptsize\sffamily\bfseries] at (0,30)
    {\iflangja{(a) タスク分割}{(a) task decomposition}};
  \foreach \i in {0,...,7} {
    \draw[fill=white] (\i*5,20) rectangle ++(5,5) node[midway] {\i};
  }
  \node[anchor=east, inner sep=1pt] at (0,22.5) {$A$};
  \node[core] (c1) at (12,8) {\iflangja{コア 1}{Core 1}};
  \node[core] (c2) at (32,8) {\iflangja{コア 2}{Core 2}};
  \draw[->] (10,20) -- (c1);
  \draw[->] (30,20) -- (c2);
  \node[res, below=2mm of c1] {\iflangja{合計}{sum}};
  \node[res, below=2mm of c2] {\iflangja{最小・最大}{min, max}};
  % (b) data decomposition
  \node[anchor=west, font=\scriptsize\sffamily\bfseries] at (58,30)
    {\iflangja{(b) データ分割}{(b) data decomposition}};
  \foreach \i in {0,...,3} {
    \draw[fill=ln-accent!10] (58+\i*5,20) rectangle ++(5,5) node[midway] {\i};
  }
  \foreach \i in {4,...,7} {
    \draw[fill=ln-tip!12] (60+\i*5,20) rectangle ++(5,5) node[midway] {\i};
  }
  \node[anchor=east, inner sep=1pt] at (58,22.5) {$A$};
  \node[core] (d1) at (68,8) {\iflangja{コア 1}{Core 1}};
  \node[core] (d2) at (92,8) {\iflangja{コア 2}{Core 2}};
  \draw[->] (68,20) -- (d1);
  \draw[->] (92,20) -- (d2);
  \node[res, below=2mm of d1] (p1) {\iflangja{部分結果 1}{partial 1}};
  \node[res, below=2mm of d2] (p2) {\iflangja{部分結果 2}{partial 2}};
  \node[res, fill=ln-source!12, minimum width=20mm] (tot) at (80,-14)
    {\iflangja{リダクション：合計，最小，最大}{reduction: sum, min, max}};
  \draw[->] (p1) -- (tot);
  \draw[->] (p2) -- (tot);
\end{tikzpicture}

% Width: body (120 mm); the drawing is about 118 mm wide.
\begin{tikzpicture}[x=1mm, y=1mm, font=\scriptsize\sffamily,
    >={Stealth[length=1.6mm]},
    core/.style={draw, rounded corners=1pt, fill=ln-box, minimum width=12mm,
                 minimum height=5mm},
    res/.style={draw, fill=ln-accent!12, minimum width=16mm, minimum height=5mm,
                align=center}]
  % (a) task decomposition
  \node[anchor=west, font=\scriptsize\sffamily\bfseries] at (0,30)
    {\iflangja{(a) タスク分割}{(a) task decomposition}};
  \foreach \i in {0,...,7} {
    \draw[fill=white] (\i*5,20) rectangle ++(5,5) node[midway] {\i};
  }
  \node[anchor=east, inner sep=1pt] at (0,22.5) {$A$};
  \node[core] (c1) at (12,8) {\iflangja{コア 1}{Core 1}};
  \node[core] (c2) at (32,8) {\iflangja{コア 2}{Core 2}};
  \draw[->] (10,20) -- (c1);
  \draw[->] (30,20) -- (c2);
  \node[res, below=2mm of c1] {\iflangja{合計}{sum}};
  \node[res, below=2mm of c2] {\iflangja{最小・最大}{min, max}};
  % (b) data decomposition
  \node[anchor=west, font=\scriptsize\sffamily\bfseries] at (58,30)
    {\iflangja{(b) データ分割}{(b) data decomposition}};
  \foreach \i in {0,...,3} {
    \draw[fill=ln-accent!10] (58+\i*5,20) rectangle ++(5,5) node[midway] {\i};
  }
  \foreach \i in {4,...,7} {
    \draw[fill=ln-tip!12] (60+\i*5,20) rectangle ++(5,5) node[midway] {\i};
  }
  \node[anchor=east, inner sep=1pt] at (58,22.5) {$A$};
  \node[core] (d1) at (68,8) {\iflangja{コア 1}{Core 1}};
  \node[core] (d2) at (92,8) {\iflangja{コア 2}{Core 2}};
  \draw[->] (68,20) -- (d1);
  \draw[->] (92,20) -- (d2);
  \node[res, below=2mm of d1] (p1) {\iflangja{部分結果 1}{partial 1}};
  \node[res, below=2mm of d2] (p2) {\iflangja{部分結果 2}{partial 2}};
  \node[res, fill=ln-source!12, minimum width=20mm] (tot) at (80,-14)
    {\iflangja{リダクション：合計，最小，最大}{reduction: sum, min, max}};
  \draw[->] (p1) -- (tot);
  \draw[->] (p2) -- (tot);
\end{tikzpicture}

% Width: body (120 mm); the drawing is about 118 mm wide.
\begin{tikzpicture}[x=1mm, y=1mm, font=\scriptsize\sffamily,
    >={Stealth[length=1.6mm]},
    core/.style={draw, rounded corners=1pt, fill=ln-box, minimum width=12mm,
                 minimum height=5mm},
    res/.style={draw, fill=ln-accent!12, minimum width=16mm, minimum height=5mm,
                align=center}]
  % (a) task decomposition
  \node[anchor=west, font=\scriptsize\sffamily\bfseries] at (0,30)
    {\iflangja{(a) タスク分割}{(a) task decomposition}};
  \foreach \i in {0,...,7} {
    \draw[fill=white] (\i*5,20) rectangle ++(5,5) node[midway] {\i};
  }
  \node[anchor=east, inner sep=1pt] at (0,22.5) {$A$};
  \node[core] (c1) at (12,8) {\iflangja{コア 1}{Core 1}};
  \node[core] (c2) at (32,8) {\iflangja{コア 2}{Core 2}};
  \draw[->] (10,20) -- (c1);
  \draw[->] (30,20) -- (c2);
  \node[res, below=2mm of c1] {\iflangja{合計}{sum}};
  \node[res, below=2mm of c2] {\iflangja{最小・最大}{min, max}};
  % (b) data decomposition
  \node[anchor=west, font=\scriptsize\sffamily\bfseries] at (58,30)
    {\iflangja{(b) データ分割}{(b) data decomposition}};
  \foreach \i in {0,...,3} {
    \draw[fill=ln-accent!10] (58+\i*5,20) rectangle ++(5,5) node[midway] {\i};
  }
  \foreach \i in {4,...,7} {
    \draw[fill=ln-tip!12] (60+\i*5,20) rectangle ++(5,5) node[midway] {\i};
  }
  \node[anchor=east, inner sep=1pt] at (58,22.5) {$A$};
  \node[core] (d1) at (68,8) {\iflangja{コア 1}{Core 1}};
  \node[core] (d2) at (92,8) {\iflangja{コア 2}{Core 2}};
  \draw[->] (68,20) -- (d1);
  \draw[->] (92,20) -- (d2);
  \node[res, below=2mm of d1] (p1) {\iflangja{部分結果 1}{partial 1}};
  \node[res, below=2mm of d2] (p2) {\iflangja{部分結果 2}{partial 2}};
  \node[res, fill=ln-source!12, minimum width=20mm] (tot) at (80,-14)
    {\iflangja{リダクション：合計，最小，最大}{reduction: sum, min, max}};
  \draw[->] (p1) -- (tot);
  \draw[->] (p2) -- (tot);
\end{tikzpicture}

  \caption{Task decomposition sends the whole array to each core and
    gives each a different operation; data decomposition gives each core
    part of the array and combines the partial results with a reduction.}
  \label{fig:l01-decomposition}
\end{figure}

\subsection{Flynn's taxonomy}

\begin{marginfigure}
  \resizebox{\linewidth}{!}{% Flynn's taxonomy as a 2x2 grid (Lesson 1).
% Rows: single or multiple instruction streams; columns: single or
% multiple data streams.  Redraws the outline figure m1l3_flynn.
% Usage: % Flynn's taxonomy as a 2x2 grid (Lesson 1).
% Rows: single or multiple instruction streams; columns: single or
% multiple data streams.  Redraws the outline figure m1l3_flynn.
% Usage: % Flynn's taxonomy as a 2x2 grid (Lesson 1).
% Rows: single or multiple instruction streams; columns: single or
% multiple data streams.  Redraws the outline figure m1l3_flynn.
% Usage: \input{figures/l01-flynn}  (inside marginfigure, \resizebox to
% \linewidth)
% Width: margin (50 mm); the drawing is about 50 mm wide.
\begin{tikzpicture}[x=1mm, y=1mm, font=\scriptsize\sffamily,
    cell/.style={draw, minimum width=19mm, minimum height=9mm, align=center,
                 fill=ln-box, inner sep=1pt}]
  \node[cell] at (19,12) {SISD};
  \node[cell, fill=ln-accent!15] at (39,12) {SIMD / SIMT};
  \node[cell] at (19,2) {MISD};
  \node[cell] at (39,2) {MIMD};
  \node[anchor=south, align=center, text=ln-muted] at (19,17)
    {\iflangja{単一データ}{single data}};
  \node[anchor=south, align=center, text=ln-muted] at (39,17)
    {\iflangja{複数データ}{multiple data}};
  \node[anchor=east, align=right, text=ln-muted] at (9,12)
    {\iflangja{単一\\命令}{single\\instruction}};
  \node[anchor=east, align=right, text=ln-muted] at (9,2)
    {\iflangja{複数\\命令}{multiple\\instruction}};
\end{tikzpicture}
  (inside marginfigure, \resizebox to
% \linewidth)
% Width: margin (50 mm); the drawing is about 50 mm wide.
\begin{tikzpicture}[x=1mm, y=1mm, font=\scriptsize\sffamily,
    cell/.style={draw, minimum width=19mm, minimum height=9mm, align=center,
                 fill=ln-box, inner sep=1pt}]
  \node[cell] at (19,12) {SISD};
  \node[cell, fill=ln-accent!15] at (39,12) {SIMD / SIMT};
  \node[cell] at (19,2) {MISD};
  \node[cell] at (39,2) {MIMD};
  \node[anchor=south, align=center, text=ln-muted] at (19,17)
    {\iflangja{単一データ}{single data}};
  \node[anchor=south, align=center, text=ln-muted] at (39,17)
    {\iflangja{複数データ}{multiple data}};
  \node[anchor=east, align=right, text=ln-muted] at (9,12)
    {\iflangja{単一\\命令}{single\\instruction}};
  \node[anchor=east, align=right, text=ln-muted] at (9,2)
    {\iflangja{複数\\命令}{multiple\\instruction}};
\end{tikzpicture}
  (inside marginfigure, \resizebox to
% \linewidth)
% Width: margin (50 mm); the drawing is about 50 mm wide.
\begin{tikzpicture}[x=1mm, y=1mm, font=\scriptsize\sffamily,
    cell/.style={draw, minimum width=19mm, minimum height=9mm, align=center,
                 fill=ln-box, inner sep=1pt}]
  \node[cell] at (19,12) {SISD};
  \node[cell, fill=ln-accent!15] at (39,12) {SIMD / SIMT};
  \node[cell] at (19,2) {MISD};
  \node[cell] at (39,2) {MIMD};
  \node[anchor=south, align=center, text=ln-muted] at (19,17)
    {\iflangja{単一データ}{single data}};
  \node[anchor=south, align=center, text=ln-muted] at (39,17)
    {\iflangja{複数データ}{multiple data}};
  \node[anchor=east, align=right, text=ln-muted] at (9,12)
    {\iflangja{単一\\命令}{single\\instruction}};
  \node[anchor=east, align=right, text=ln-muted] at (9,2)
    {\iflangja{複数\\命令}{multiple\\instruction}};
\end{tikzpicture}
}
  \caption{Flynn's taxonomy.  GPUs sit in the single-instruction,
    multiple-data quadrant as SIMT.}
  \label{fig:l01-flynn}
\end{marginfigure}

\term{Flynn's taxonomy}[フリンの分類] classifies a machine by the number
of instruction streams and the number of data streams it processes at a
time (\cref{fig:l01-flynn}): SISD, SIMD, MISD (multiple instruction,
single data, a quadrant with almost no commercial examples) and
\term{MIMD}[MIMD] (multiple instruction, multiple data), in which
independent processors run independent instruction streams on
independent data, as in a multicore CPU.  A CPU can be SISD (one core, scalar code),
SIMD (its vector units) or MIMD (its many cores).  A GPU is placed in
the SIMD quadrant under the name \term{SIMT}[SIMT] (single instruction,
multiple threads): like SIMD, one fetched instruction drives many
operations, but the operations belong to distinct threads, each with its
own registers, executed on many ALUs rather than on a vector unit; the
group shares one program counter and an active mask
(\cref{sec:l01-gpu-architecture}, Lesson~6).  The distinction is
developed in \cref{sec:l01-gpu-architecture} and in
Lesson~4.\caution{CPG §4.1 SIMT Architecture draws the line as: SIMD
exposes the vector width to software, SIMT does not; a SIMT thread is
written as an ordinary scalar thread.  Since Volta the hardware also
keeps a program counter per thread (Independent Thread Scheduling)
\cite{nvidia2024cudaguide}; Lesson~5 returns to it.}

\subsection{SPMD on SIMT hardware}

In the SPMD style all threads run the same program and differ only in
the data they touch; data decomposition is its natural pattern.  SPMD is
the typical way of programming a GPU, and SIMT is the execution model of
the hardware that runs it: the programmer writes one scalar program for
one thread, the hardware fetches each of its instructions once for a
whole group of threads.

A CUDA program consists of host code and a kernel, the function that
runs on the GPU.  Every thread executes the same kernel and uses its
thread index to select its own data.  The kernel below squares one
element per thread; the index arithmetic and the launch syntax are
explained in Lesson~3.
\Needspace{8\baselineskip}
\begin{lstlisting}[language=C++, numbers=none]
__global__ void square(const float *A, float *B, int n)
{
  int i = blockIdx.x * blockDim.x + threadIdx.x;
  if (i < n) B[i] = A[i] * A[i];
}
\end{lstlisting}
Summing the array in the same style is not as simple: a kernel in which
every thread executes \texttt{sum += A[i]} on one shared variable is a
race, since thousands of threads read, add and write the same location
concurrently.  The sum needs a reduction, built from per-thread partial
sums and a combining step (Lesson~2 for OpenMP, Lesson~5 for the GPU
primitives).\caution{\texttt{sum += A[i]} from every thread of a kernel
is incorrect, not merely slow: the result depends on the interleaving.
Treat any shared accumulator as a reduction.}

\section{Introduction to GPU architecture}
\label{sec:l01-gpu-architecture}

A GPU is a multiprocessor of many cores, each of which runs many threads.
This section names the parts of such a core, sketches its pipeline and
defines the unit in which it executes threads.\source{M1, GPU
architecture overview and warps; K12 §1.2.1; A18 ch.~1; CPG §4.1 SIMT
Architecture.}

\subsection{Streaming multiprocessors and lanes}

NVIDIA calls the core of a GPU a
\term{streaming multiprocessor (SM)}[ストリーミング・マルチプロセッサ]\index{SM|see{streaming multiprocessor (SM)}},
also a SIMT multiprocessor.\margin{K12 §1.2.1: NVIDIA calls the core an
SM and AMD a compute unit (CU); each holds a number of SIMD units, also
called stream processors \cite{kim2012gpgpu}.}  Inside an SM
(\cref{fig:l01-sm}) sit an instruction cache, a decoder, a block of
shared memory, a set of program counters and register files for the
threads it holds, and many execution units.  Each execution unit is a
\term{SIMT lane}[SIMTレーン]: one ALU that executes one thread's share
of each instruction.  NVIDIA markets a lane as a
\term{stream processor (CUDA core)}[ストリーミングプロセッサ（CUDAコア）]\index{CUDA core|see{stream processor (CUDA core)}},
which invites confusion with a CPU core.\caution{In this course
\emph{core} $\ne$ stream processor.  A core is the SM; a stream
processor is one lane of it.  A \enquote{\num{16896}-core} H100 is a
132-SM GPU with 128 lanes per SM.}  This book therefore fixes the
vocabulary: a \emph{core} is an SM, and a SIMT lane is never called a
core.

\begin{figure}[htbp]
  \centering
  % A GPU as a grid of streaming multiprocessors, one SM zoomed (Lesson 1).
% Left: 4x4 "core" (SM) grid with a cache bar.  Right: inside one SM,
% a stack of PCs, the instruction cache, decoder, shared memory, a row of
% stream processors (SIMT lanes) and the caches.
% Redraws the outline figure m1l4_gpuarch.
% Usage: % A GPU as a grid of streaming multiprocessors, one SM zoomed (Lesson 1).
% Left: 4x4 "core" (SM) grid with a cache bar.  Right: inside one SM,
% a stack of PCs, the instruction cache, decoder, shared memory, a row of
% stream processors (SIMT lanes) and the caches.
% Redraws the outline figure m1l4_gpuarch.
% Usage: % A GPU as a grid of streaming multiprocessors, one SM zoomed (Lesson 1).
% Left: 4x4 "core" (SM) grid with a cache bar.  Right: inside one SM,
% a stack of PCs, the instruction cache, decoder, shared memory, a row of
% stream processors (SIMT lanes) and the caches.
% Redraws the outline figure m1l4_gpuarch.
% Usage: \input{figures/l01-sm}
% Width: body (120 mm); the drawing is about 118 mm wide.
\begin{tikzpicture}[x=1mm, y=1mm, font=\scriptsize\sffamily,
    >={Stealth[length=1.6mm]},
    blk/.style={draw, fill=ln-box, minimum height=5mm, inner sep=2pt,
                align=center}]
  % GPU grid
  \node[font=\small\sffamily\bfseries] at (20,37) {GPU};
  \foreach \i in {0,...,3} \foreach \j in {0,1} {
    \draw[fill=ln-box] (\i*10,\j*6+22) rectangle ++(9,5)
      node[midway] {SM};
  }
  \draw[fill=ln-accent!15] (0,18) rectangle (39,21)
    node[midway] {\iflangja{キャッシュ}{Cache}};
  \foreach \i in {0,...,3} \foreach \j in {0,1} {
    \draw[fill=ln-box] (\i*10,\j*6+6) rectangle ++(9,5)
      node[midway] {SM};
  }
  \draw[fill=ln-accent!15] (0,3) rectangle (39,5);
  % zoom lines
  \draw[ln-muted, densely dashed] (30,27) -- (52,36);
  \draw[ln-muted, densely dashed] (30,22) -- (52,0);
  \draw[fill=ln-box, thick, ln-accent] (30,22) rectangle (39,27)
    node[midway, text=black] {SM};
  % inside one SM
  \draw[ln-accent, thick, rounded corners=1pt] (52,0) rectangle (118,36);
  \node[anchor=north west, font=\scriptsize\sffamily\bfseries, text=ln-accent]
    at (53,36) {\iflangja{1個の SM}{one SM}};
  \foreach \k in {0,1,2} {
    \draw[fill=white] (55+\k*1.2,24+\k*1.2) rectangle ++(9,5);
  }
  \node at (61.9,28.9) {PC};
  \node[blk, anchor=west, minimum width=14mm] at (68,27) {\iflangja{命令\\キャッシュ}{I-cache}};
  \node[blk, anchor=west, minimum width=14mm] at (84,27) {\iflangja{デコーダ}{Decoder}};
  \node[blk, anchor=west, minimum width=16mm] at (100,27) {\iflangja{共有\\メモリ}{Shared\\memory}};
  \foreach \k in {0,1,2} {
    \draw[fill=ln-accent!15] (55+\k*9,8) rectangle ++(8,12)
      node[midway, align=center, rotate=90] {SP};
  }
  \node at (85,14) {$\cdots$};
  \draw[fill=ln-accent!15] (90,8) rectangle ++(8,12)
    node[midway, rotate=90] {SP};
  \node[anchor=west, align=left, text=ln-muted] at (99,14)
    {\iflangja{ストリーミング\\プロセッサ\\＝SIMT レーン}{stream\\processors\\= SIMT lanes}};
  \node[blk, anchor=west, minimum width=61mm] at (55,4) {\iflangja{キャッシュ}{Caches}};
\end{tikzpicture}

% Width: body (120 mm); the drawing is about 118 mm wide.
\begin{tikzpicture}[x=1mm, y=1mm, font=\scriptsize\sffamily,
    >={Stealth[length=1.6mm]},
    blk/.style={draw, fill=ln-box, minimum height=5mm, inner sep=2pt,
                align=center}]
  % GPU grid
  \node[font=\small\sffamily\bfseries] at (20,37) {GPU};
  \foreach \i in {0,...,3} \foreach \j in {0,1} {
    \draw[fill=ln-box] (\i*10,\j*6+22) rectangle ++(9,5)
      node[midway] {SM};
  }
  \draw[fill=ln-accent!15] (0,18) rectangle (39,21)
    node[midway] {\iflangja{キャッシュ}{Cache}};
  \foreach \i in {0,...,3} \foreach \j in {0,1} {
    \draw[fill=ln-box] (\i*10,\j*6+6) rectangle ++(9,5)
      node[midway] {SM};
  }
  \draw[fill=ln-accent!15] (0,3) rectangle (39,5);
  % zoom lines
  \draw[ln-muted, densely dashed] (30,27) -- (52,36);
  \draw[ln-muted, densely dashed] (30,22) -- (52,0);
  \draw[fill=ln-box, thick, ln-accent] (30,22) rectangle (39,27)
    node[midway, text=black] {SM};
  % inside one SM
  \draw[ln-accent, thick, rounded corners=1pt] (52,0) rectangle (118,36);
  \node[anchor=north west, font=\scriptsize\sffamily\bfseries, text=ln-accent]
    at (53,36) {\iflangja{1個の SM}{one SM}};
  \foreach \k in {0,1,2} {
    \draw[fill=white] (55+\k*1.2,24+\k*1.2) rectangle ++(9,5);
  }
  \node at (61.9,28.9) {PC};
  \node[blk, anchor=west, minimum width=14mm] at (68,27) {\iflangja{命令\\キャッシュ}{I-cache}};
  \node[blk, anchor=west, minimum width=14mm] at (84,27) {\iflangja{デコーダ}{Decoder}};
  \node[blk, anchor=west, minimum width=16mm] at (100,27) {\iflangja{共有\\メモリ}{Shared\\memory}};
  \foreach \k in {0,1,2} {
    \draw[fill=ln-accent!15] (55+\k*9,8) rectangle ++(8,12)
      node[midway, align=center, rotate=90] {SP};
  }
  \node at (85,14) {$\cdots$};
  \draw[fill=ln-accent!15] (90,8) rectangle ++(8,12)
    node[midway, rotate=90] {SP};
  \node[anchor=west, align=left, text=ln-muted] at (99,14)
    {\iflangja{ストリーミング\\プロセッサ\\＝SIMT レーン}{stream\\processors\\= SIMT lanes}};
  \node[blk, anchor=west, minimum width=61mm] at (55,4) {\iflangja{キャッシュ}{Caches}};
\end{tikzpicture}

% Width: body (120 mm); the drawing is about 118 mm wide.
\begin{tikzpicture}[x=1mm, y=1mm, font=\scriptsize\sffamily,
    >={Stealth[length=1.6mm]},
    blk/.style={draw, fill=ln-box, minimum height=5mm, inner sep=2pt,
                align=center}]
  % GPU grid
  \node[font=\small\sffamily\bfseries] at (20,37) {GPU};
  \foreach \i in {0,...,3} \foreach \j in {0,1} {
    \draw[fill=ln-box] (\i*10,\j*6+22) rectangle ++(9,5)
      node[midway] {SM};
  }
  \draw[fill=ln-accent!15] (0,18) rectangle (39,21)
    node[midway] {\iflangja{キャッシュ}{Cache}};
  \foreach \i in {0,...,3} \foreach \j in {0,1} {
    \draw[fill=ln-box] (\i*10,\j*6+6) rectangle ++(9,5)
      node[midway] {SM};
  }
  \draw[fill=ln-accent!15] (0,3) rectangle (39,5);
  % zoom lines
  \draw[ln-muted, densely dashed] (30,27) -- (52,36);
  \draw[ln-muted, densely dashed] (30,22) -- (52,0);
  \draw[fill=ln-box, thick, ln-accent] (30,22) rectangle (39,27)
    node[midway, text=black] {SM};
  % inside one SM
  \draw[ln-accent, thick, rounded corners=1pt] (52,0) rectangle (118,36);
  \node[anchor=north west, font=\scriptsize\sffamily\bfseries, text=ln-accent]
    at (53,36) {\iflangja{1個の SM}{one SM}};
  \foreach \k in {0,1,2} {
    \draw[fill=white] (55+\k*1.2,24+\k*1.2) rectangle ++(9,5);
  }
  \node at (61.9,28.9) {PC};
  \node[blk, anchor=west, minimum width=14mm] at (68,27) {\iflangja{命令\\キャッシュ}{I-cache}};
  \node[blk, anchor=west, minimum width=14mm] at (84,27) {\iflangja{デコーダ}{Decoder}};
  \node[blk, anchor=west, minimum width=16mm] at (100,27) {\iflangja{共有\\メモリ}{Shared\\memory}};
  \foreach \k in {0,1,2} {
    \draw[fill=ln-accent!15] (55+\k*9,8) rectangle ++(8,12)
      node[midway, align=center, rotate=90] {SP};
  }
  \node at (85,14) {$\cdots$};
  \draw[fill=ln-accent!15] (90,8) rectangle ++(8,12)
    node[midway, rotate=90] {SP};
  \node[anchor=west, align=left, text=ln-muted] at (99,14)
    {\iflangja{ストリーミング\\プロセッサ\\＝SIMT レーン}{stream\\processors\\= SIMT lanes}};
  \node[blk, anchor=west, minimum width=61mm] at (55,4) {\iflangja{キャッシュ}{Caches}};
\end{tikzpicture}

  \caption{A GPU as a grid of streaming multiprocessors sharing a cache,
    and the contents of one SM: program counters, instruction cache,
    decoder, shared memory, stream processors (SIMT lanes) and caches.}
  \label{fig:l01-sm}
\end{figure}

\subsection{The GPU pipeline}

The pipeline of an SM (\cref{fig:l01-gpu-pipeline}) has the stages of
the CPU pipeline of \cref{sec:l01-processor-paradigms}, arranged for
many threads.  The fetch stage keeps one program counter per warp, the
thread group defined below, and fetches for whichever warp a scheduler
selects,\margin{The word \emph{warp} comes from weaving, \enquote{the
first parallel thread technology}: CPG §4.1 SIMT Architecture
\cite{nvidia2024cudaguide}.} by a
round-robin policy or by a greedy policy informed by cache misses and
branch outcomes.  Decode and register read follow.  A second scheduler
with a scoreboard then picks a group whose source operands are ready,
the lanes execute the instruction, a memory-access unit handles loads
and stores, and the write-back stage returns results to the
registers.\margin{Lesson~4 traces this pipeline cycle by cycle with its
scoreboard and register banks; Lesson~7 compares the scheduling
policies.}

\begin{figure}[htbp]
  \centering
  % GPU pipeline overview (Lesson 1).
% Fetch (from the I-cache, with one PC per warp and an active mask) ->
% decode -> register read -> scheduler with a scoreboard -> SIMT lanes ->
% memory access -> write-back, with the result fed back to the registers.
% Redraws the outline figure m1l4_gpupipeline.
% Usage: % GPU pipeline overview (Lesson 1).
% Fetch (from the I-cache, with one PC per warp and an active mask) ->
% decode -> register read -> scheduler with a scoreboard -> SIMT lanes ->
% memory access -> write-back, with the result fed back to the registers.
% Redraws the outline figure m1l4_gpupipeline.
% Usage: % GPU pipeline overview (Lesson 1).
% Fetch (from the I-cache, with one PC per warp and an active mask) ->
% decode -> register read -> scheduler with a scoreboard -> SIMT lanes ->
% memory access -> write-back, with the result fed back to the registers.
% Redraws the outline figure m1l4_gpupipeline.
% Usage: \input{figures/l01-gpu-pipeline}
% Width: body (120 mm); the drawing is about 118 mm wide.
\begin{tikzpicture}[x=1mm, y=1mm, font=\scriptsize\sffamily,
    >={Stealth[length=1.6mm]},
    st/.style={draw, fill=ln-box, minimum height=7mm, minimum width=15mm,
               align=center, inner sep=2pt},
    mem/.style={draw, fill=ln-accent!15, minimum height=6mm, align=center,
                inner sep=2pt}]
  \node[st] (fetch) at (10,20) {\iflangja{フェッチ}{Fetch}};
  \node[st, right=4mm of fetch] (dec) {\iflangja{デコード}{Decode}};
  \node[st, right=4mm of dec] (reg) {\iflangja{レジスタ\\読み出し}{Register\\read}};
  \node[st, right=4mm of reg, minimum width=19mm] (sched)
    {\iflangja{スケジューラ}{Scheduler}};
  \node[st, right=4mm of sched, minimum width=19mm, fill=ln-accent!12] (lanes)
    {\iflangja{レーン $0 \ldots N\!-\!1$}{Lanes $0 \ldots N\!-\!1$}};
  \draw[->] (fetch) -- (dec);
  \draw[->] (dec) -- (reg);
  \draw[->] (reg) -- (sched);
  \draw[->] (sched) -- (lanes);
  % memory access and write-back on a second row
  \node[st, below=5mm of lanes, minimum width=19mm] (ma)
    {\iflangja{メモリアクセス}{Memory access}};
  \node[st, left=4mm of ma, minimum width=15mm] (wb)
    {\iflangja{ライトバック}{Write-back}};
  \draw[->] (lanes) -- (ma);
  \draw[->] (ma) -- (wb);
  \draw[->] (wb.west) -| (reg.south);
  % per-warp state above fetch, I-cache below it, scoreboard above the scheduler
  \node[mem, above=4mm of fetch, minimum width=17mm] (pcs)
    {\iflangja{ワープごとの PC\\アクティブマスク}{per-warp PC,\\active mask}};
  \draw[<->] (pcs) -- (fetch);
  \node[mem, below=4mm of fetch, minimum width=15mm] (icache) {\iflangja{命令キャッシュ}{I-cache}};
  \draw[<->] (icache) -- (fetch);
  \node[mem, above=4mm of sched, minimum width=19mm] (sb)
    {\iflangja{スコアボード}{scoreboard}};
  \draw[<->] (sb) -- (sched);
  \node[mem, right=3mm of ma, minimum width=10mm] (dmem) {\iflangja{メモリ}{Memory}};
  \draw[<->] (ma) -- (dmem);
\end{tikzpicture}

% Width: body (120 mm); the drawing is about 118 mm wide.
\begin{tikzpicture}[x=1mm, y=1mm, font=\scriptsize\sffamily,
    >={Stealth[length=1.6mm]},
    st/.style={draw, fill=ln-box, minimum height=7mm, minimum width=15mm,
               align=center, inner sep=2pt},
    mem/.style={draw, fill=ln-accent!15, minimum height=6mm, align=center,
                inner sep=2pt}]
  \node[st] (fetch) at (10,20) {\iflangja{フェッチ}{Fetch}};
  \node[st, right=4mm of fetch] (dec) {\iflangja{デコード}{Decode}};
  \node[st, right=4mm of dec] (reg) {\iflangja{レジスタ\\読み出し}{Register\\read}};
  \node[st, right=4mm of reg, minimum width=19mm] (sched)
    {\iflangja{スケジューラ}{Scheduler}};
  \node[st, right=4mm of sched, minimum width=19mm, fill=ln-accent!12] (lanes)
    {\iflangja{レーン $0 \ldots N\!-\!1$}{Lanes $0 \ldots N\!-\!1$}};
  \draw[->] (fetch) -- (dec);
  \draw[->] (dec) -- (reg);
  \draw[->] (reg) -- (sched);
  \draw[->] (sched) -- (lanes);
  % memory access and write-back on a second row
  \node[st, below=5mm of lanes, minimum width=19mm] (ma)
    {\iflangja{メモリアクセス}{Memory access}};
  \node[st, left=4mm of ma, minimum width=15mm] (wb)
    {\iflangja{ライトバック}{Write-back}};
  \draw[->] (lanes) -- (ma);
  \draw[->] (ma) -- (wb);
  \draw[->] (wb.west) -| (reg.south);
  % per-warp state above fetch, I-cache below it, scoreboard above the scheduler
  \node[mem, above=4mm of fetch, minimum width=17mm] (pcs)
    {\iflangja{ワープごとの PC\\アクティブマスク}{per-warp PC,\\active mask}};
  \draw[<->] (pcs) -- (fetch);
  \node[mem, below=4mm of fetch, minimum width=15mm] (icache) {\iflangja{命令キャッシュ}{I-cache}};
  \draw[<->] (icache) -- (fetch);
  \node[mem, above=4mm of sched, minimum width=19mm] (sb)
    {\iflangja{スコアボード}{scoreboard}};
  \draw[<->] (sb) -- (sched);
  \node[mem, right=3mm of ma, minimum width=10mm] (dmem) {\iflangja{メモリ}{Memory}};
  \draw[<->] (ma) -- (dmem);
\end{tikzpicture}

% Width: body (120 mm); the drawing is about 118 mm wide.
\begin{tikzpicture}[x=1mm, y=1mm, font=\scriptsize\sffamily,
    >={Stealth[length=1.6mm]},
    st/.style={draw, fill=ln-box, minimum height=7mm, minimum width=15mm,
               align=center, inner sep=2pt},
    mem/.style={draw, fill=ln-accent!15, minimum height=6mm, align=center,
                inner sep=2pt}]
  \node[st] (fetch) at (10,20) {\iflangja{フェッチ}{Fetch}};
  \node[st, right=4mm of fetch] (dec) {\iflangja{デコード}{Decode}};
  \node[st, right=4mm of dec] (reg) {\iflangja{レジスタ\\読み出し}{Register\\read}};
  \node[st, right=4mm of reg, minimum width=19mm] (sched)
    {\iflangja{スケジューラ}{Scheduler}};
  \node[st, right=4mm of sched, minimum width=19mm, fill=ln-accent!12] (lanes)
    {\iflangja{レーン $0 \ldots N\!-\!1$}{Lanes $0 \ldots N\!-\!1$}};
  \draw[->] (fetch) -- (dec);
  \draw[->] (dec) -- (reg);
  \draw[->] (reg) -- (sched);
  \draw[->] (sched) -- (lanes);
  % memory access and write-back on a second row
  \node[st, below=5mm of lanes, minimum width=19mm] (ma)
    {\iflangja{メモリアクセス}{Memory access}};
  \node[st, left=4mm of ma, minimum width=15mm] (wb)
    {\iflangja{ライトバック}{Write-back}};
  \draw[->] (lanes) -- (ma);
  \draw[->] (ma) -- (wb);
  \draw[->] (wb.west) -| (reg.south);
  % per-warp state above fetch, I-cache below it, scoreboard above the scheduler
  \node[mem, above=4mm of fetch, minimum width=17mm] (pcs)
    {\iflangja{ワープごとの PC\\アクティブマスク}{per-warp PC,\\active mask}};
  \draw[<->] (pcs) -- (fetch);
  \node[mem, below=4mm of fetch, minimum width=15mm] (icache) {\iflangja{命令キャッシュ}{I-cache}};
  \draw[<->] (icache) -- (fetch);
  \node[mem, above=4mm of sched, minimum width=19mm] (sb)
    {\iflangja{スコアボード}{scoreboard}};
  \draw[<->] (sb) -- (sched);
  \node[mem, right=3mm of ma, minimum width=10mm] (dmem) {\iflangja{メモリ}{Memory}};
  \draw[<->] (ma) -- (dmem);
\end{tikzpicture}

  \caption{The GPU pipeline: fetch with one program counter per warp,
    decode, register read, a scheduler with a scoreboard that selects a
    ready warp, the SIMT lanes, memory access and write-back.}
  \label{fig:l01-gpu-pipeline}
\end{figure}

\subsection{Warps}

\begin{definition}[Warp]
  \label{def:l01-warp}
  A \term{warp}[ワープ] is a group of threads that the hardware executes
  as a unit: one instruction is fetched per warp and executed for all
  of its threads, one thread per lane.  The number of threads in the
  group is the \term{warp size}[ワープサイズ].  AMD calls the same unit
  a \term{wavefront}[ウェーブフロント]; this book uses the two words
  interchangeably.
\end{definition}

The warp is the fundamental unit of execution on a GPU.  Its size has
stayed at 32 on NVIDIA hardware for a long time, but the architecture
does not fix it, and future hardware may change it.\margin{Warp size 32
for every compute capability from 5.0 to 9.0: CPG Table~21, Technical
Specifications per Compute Capability \cite{nvidia2024cudaguide}.}  A
kernel launch names a number of threads, grouped into thread blocks
(Lesson~3); the hardware cuts each block into warps of the warp size.  When a warp's source operands are ready, the
whole warp executes the instruction; then another ready warp is
selected, and so on.  The warp need not be as
wide as the SM: a 32-thread warp on an SM with 8 lanes issues over four
cycles.\margin{G80 (2006) had 8 lanes per SM, so a 32-thread warp took
four cycles per instruction: K12 §1.4.1 \cite{kim2012gpgpu}.  Current
SMs have 128 FP32 lanes split into four partitions
\cite{nvidia2022hopper}.}

\begin{example}
  \label{ex:l01-warps}
  A kernel of three instructions is launched with 12 threads (one
  thread block) on an SM whose warp size is 4 and which has 4 lanes (\cref{fig:l01-warp-execution}).
  The threads form warps $W_0$, $W_1$, $W_2$.\tip{Threads per block
  $\div$ warp size, rounded up, is the number of warps per block; a
  block of 100 threads with warp size 32 makes four warps, the last one
  only $4/32$ full.}  With a round-robin
  scheduler, one warp instruction per cycle and no stalls, the
  execution order is
  \begin{center}
    \begin{tabular}{lccccccccc}
      \toprule
      cycle & 1 & 2 & 3 & 4 & 5 & 6 & 7 & 8 & 9 \\
      \midrule
      warp & $W_0$ & $W_1$ & $W_2$ & $W_0$ & $W_1$ & $W_2$ & $W_0$ & $W_1$ & $W_2$ \\
      instruction & 1 & 1 & 1 & 2 & 2 & 2 & 3 & 3 & 3 \\
      \bottomrule
    \end{tabular}
  \end{center}
  Nine instruction fetches execute $12 \times 3 = 36$ thread
  instructions.  If instruction~2 of $W_0$ had to wait for the product of
  instruction~1 (a multiplier latency longer than the three cycles
  between them), the scheduler would skip $W_0$ in cycle~4 and issue $W_1$ and
  $W_2$ first: this is the multithreading of
  \cref{sec:l01-processor-paradigms} at the granularity of warps.
\end{example}

\begin{figure}[htbp]
  \centering
  % Warp execution on a SIMT unit (Lesson 1).
% A kernel of three instructions, 12 threads launched, warp size 4: the
% threads form three warps; warp 0 occupies the four lanes and executes
% instruction 1 while warps 1 and 2 wait.
% Redraws the outline figures m1l4_gpuwarp and m1l4_gpuwarp2.
% Usage: % Warp execution on a SIMT unit (Lesson 1).
% A kernel of three instructions, 12 threads launched, warp size 4: the
% threads form three warps; warp 0 occupies the four lanes and executes
% instruction 1 while warps 1 and 2 wait.
% Redraws the outline figures m1l4_gpuwarp and m1l4_gpuwarp2.
% Usage: % Warp execution on a SIMT unit (Lesson 1).
% A kernel of three instructions, 12 threads launched, warp size 4: the
% threads form three warps; warp 0 occupies the four lanes and executes
% instruction 1 while warps 1 and 2 wait.
% Redraws the outline figures m1l4_gpuwarp and m1l4_gpuwarp2.
% Usage: \input{figures/l01-warp-execution}
% Width: body (120 mm); the drawing is about 112 mm wide.
\begin{tikzpicture}[x=1mm, y=1mm, font=\scriptsize\sffamily,
    >={Stealth[length=1.6mm]}]
  % program
  \draw[fill=white] (0,8) rectangle (26,30);
  \node[anchor=north west, align=left] at (1,30)
    {\iflangja{カーネル}{kernel}\\[1pt]
     \textbf{1:} \texttt{mul r1,r2,r3}\\
     \textbf{2:} \texttt{add r4,r1,r5}\\
     \textbf{3:} \texttt{st [r6],r4}};
  \draw[ln-accent, thick, rounded corners=1pt] (0.8,23.0) rectangle (24,26.1);
  % threads: three warps of four
  \node[anchor=west] at (32,30) {\iflangja{起動された12スレッド}{12 threads launched}};
  \foreach \w/\c in {0/ln-accent!25, 1/ln-tip!25, 2/ln-source!20} {
    \foreach \t in {0,...,3} {
      \draw[fill=\c] (32+\w*20+\t*4.2,18) rectangle ++(3.6,4)
        node[midway] {T};
    }
    \draw[ln-muted] (32+\w*20,17.2) -- ++(0,-1.2) -- ++(16.2,0) -- ++(0,1.2);
    \node[below, inner sep=1pt] at (32+\w*20+8.1,16)
      {\iflangja{ワープ \w}{warp \w}};
  }
  % SIMT execution unit
  \draw[fill=ln-box] (32,0) rectangle (76,9);
  \node[anchor=north east, text=ln-muted] at (75.5,9)
    {\iflangja{実行ユニット（4レーン）}{SIMT execution unit (4 lanes)}};
  \foreach \t in {0,...,3} {
    \draw[fill=ln-accent!25] (40+\t*8,1.5) rectangle ++(6,4)
      node[midway] {SP};
  }
  \draw[->, ln-accent, thick] (48.2,15.8) -- (41,5.8)
    node[pos=0.5, right, inner sep=1pt, text=ln-accent]
    {\iflangja{ワープ 0 が命令 1 を実行}{warp 0 runs instruction 1}};
  \node[anchor=west, text=ln-muted, align=left] at (80,4)
    {\iflangja{ワープ 1，2 は\\待機}{warps 1 and 2\\wait their turn}};
\end{tikzpicture}

% Width: body (120 mm); the drawing is about 112 mm wide.
\begin{tikzpicture}[x=1mm, y=1mm, font=\scriptsize\sffamily,
    >={Stealth[length=1.6mm]}]
  % program
  \draw[fill=white] (0,8) rectangle (26,30);
  \node[anchor=north west, align=left] at (1,30)
    {\iflangja{カーネル}{kernel}\\[1pt]
     \textbf{1:} \texttt{mul r1,r2,r3}\\
     \textbf{2:} \texttt{add r4,r1,r5}\\
     \textbf{3:} \texttt{st [r6],r4}};
  \draw[ln-accent, thick, rounded corners=1pt] (0.8,23.0) rectangle (24,26.1);
  % threads: three warps of four
  \node[anchor=west] at (32,30) {\iflangja{起動された12スレッド}{12 threads launched}};
  \foreach \w/\c in {0/ln-accent!25, 1/ln-tip!25, 2/ln-source!20} {
    \foreach \t in {0,...,3} {
      \draw[fill=\c] (32+\w*20+\t*4.2,18) rectangle ++(3.6,4)
        node[midway] {T};
    }
    \draw[ln-muted] (32+\w*20,17.2) -- ++(0,-1.2) -- ++(16.2,0) -- ++(0,1.2);
    \node[below, inner sep=1pt] at (32+\w*20+8.1,16)
      {\iflangja{ワープ \w}{warp \w}};
  }
  % SIMT execution unit
  \draw[fill=ln-box] (32,0) rectangle (76,9);
  \node[anchor=north east, text=ln-muted] at (75.5,9)
    {\iflangja{実行ユニット（4レーン）}{SIMT execution unit (4 lanes)}};
  \foreach \t in {0,...,3} {
    \draw[fill=ln-accent!25] (40+\t*8,1.5) rectangle ++(6,4)
      node[midway] {SP};
  }
  \draw[->, ln-accent, thick] (48.2,15.8) -- (41,5.8)
    node[pos=0.5, right, inner sep=1pt, text=ln-accent]
    {\iflangja{ワープ 0 が命令 1 を実行}{warp 0 runs instruction 1}};
  \node[anchor=west, text=ln-muted, align=left] at (80,4)
    {\iflangja{ワープ 1，2 は\\待機}{warps 1 and 2\\wait their turn}};
\end{tikzpicture}

% Width: body (120 mm); the drawing is about 112 mm wide.
\begin{tikzpicture}[x=1mm, y=1mm, font=\scriptsize\sffamily,
    >={Stealth[length=1.6mm]}]
  % program
  \draw[fill=white] (0,8) rectangle (26,30);
  \node[anchor=north west, align=left] at (1,30)
    {\iflangja{カーネル}{kernel}\\[1pt]
     \textbf{1:} \texttt{mul r1,r2,r3}\\
     \textbf{2:} \texttt{add r4,r1,r5}\\
     \textbf{3:} \texttt{st [r6],r4}};
  \draw[ln-accent, thick, rounded corners=1pt] (0.8,23.0) rectangle (24,26.1);
  % threads: three warps of four
  \node[anchor=west] at (32,30) {\iflangja{起動された12スレッド}{12 threads launched}};
  \foreach \w/\c in {0/ln-accent!25, 1/ln-tip!25, 2/ln-source!20} {
    \foreach \t in {0,...,3} {
      \draw[fill=\c] (32+\w*20+\t*4.2,18) rectangle ++(3.6,4)
        node[midway] {T};
    }
    \draw[ln-muted] (32+\w*20,17.2) -- ++(0,-1.2) -- ++(16.2,0) -- ++(0,1.2);
    \node[below, inner sep=1pt] at (32+\w*20+8.1,16)
      {\iflangja{ワープ \w}{warp \w}};
  }
  % SIMT execution unit
  \draw[fill=ln-box] (32,0) rectangle (76,9);
  \node[anchor=north east, text=ln-muted] at (75.5,9)
    {\iflangja{実行ユニット（4レーン）}{SIMT execution unit (4 lanes)}};
  \foreach \t in {0,...,3} {
    \draw[fill=ln-accent!25] (40+\t*8,1.5) rectangle ++(6,4)
      node[midway] {SP};
  }
  \draw[->, ln-accent, thick] (48.2,15.8) -- (41,5.8)
    node[pos=0.5, right, inner sep=1pt, text=ln-accent]
    {\iflangja{ワープ 0 が命令 1 を実行}{warp 0 runs instruction 1}};
  \node[anchor=west, text=ln-muted, align=left] at (80,4)
    {\iflangja{ワープ 1，2 は\\待機}{warps 1 and 2\\wait their turn}};
\end{tikzpicture}

  \caption{Twelve launched threads form three warps of four; warp~0
    occupies the four lanes of the SIMT execution unit for instruction~1
    while the other warps wait.}
  \label{fig:l01-warp-execution}
\end{figure}

\begin{keypoint}
  A GPU is a multiprocessor of SMs; an SM is a multithreaded SIMT core
  of many lanes; a warp is the group of threads (32 on NVIDIA GPUs) that
  shares one instruction fetch and is scheduled as a unit.  All
  later performance discussion, from occupancy to divergence and
  coalescing, is in terms of warps.
\end{keypoint}

\section{History trend of GPU architectures}
\label{sec:l01-history}

GPU architecture has passed through four stages.\source{M1, graphics
pipeline and GPU evolution; K12 §1.1; A18 ch.~1; HWP \cite{nvidia2022hopper}.}
\begin{enumerate}
  \item \emph{Traditional pipelined GPUs}, with separate fixed-function
    units for vertex processing and for pixel processing;
  \item \emph{programmable GPUs}, in which the vertex and pixel units
    became programmable shaders;
  \item the progression of the shaders into unified, CUDA-programmable
    cores, which opened general-purpose computing (GPGPU);
  \item \emph{modern GPUs}, which add machine-learning accelerator
    components such as tensor cores alongside the programmable cores.
\end{enumerate}

\subsection{The traditional 3D graphics pipeline}

The \term{graphics pipeline}[グラフィックスパイプライン] of the early
GPU (\cref{fig:l01-graphics-pipeline}) takes 3D primitives, typically
triangles given by their vertices, as input.  A geometry stage transforms
and lights the vertices and projects them into 2D primitives on the
screen plane; a rasterisation stage converts each 2D primitive into the
pixels it covers and computes their colours.  The program, or originally
the fixed function, applied to each vertex is the
\term{vertex shader}[頂点シェーダ]; the one applied to each pixel (or
fragment) is the \term{pixel shader}[ピクセルシェーダ].\margin{Both
stages apply the same operation to a large number of independent
vertices or pixels: the data parallelism that GPGPU later exploited was
built into the graphics workload from the start.}

\begin{figure}[htbp]
  \centering
  % Traditional 3D graphics pipeline (Lesson 1).
% 3D primitives -> geometry stage (vertex shader) -> 2D primitives ->
% rasterisation stage (pixel/fragment shader) -> pixels.
% Redraws the outline figure m1l5_trad3dpipeline.
% Usage: % Traditional 3D graphics pipeline (Lesson 1).
% 3D primitives -> geometry stage (vertex shader) -> 2D primitives ->
% rasterisation stage (pixel/fragment shader) -> pixels.
% Redraws the outline figure m1l5_trad3dpipeline.
% Usage: % Traditional 3D graphics pipeline (Lesson 1).
% 3D primitives -> geometry stage (vertex shader) -> 2D primitives ->
% rasterisation stage (pixel/fragment shader) -> pixels.
% Redraws the outline figure m1l5_trad3dpipeline.
% Usage: \input{figures/l01-graphics-pipeline}
% Width: body (120 mm); the drawing is about 116 mm wide.
\begin{tikzpicture}[x=1mm, y=1mm, font=\scriptsize\sffamily,
    >={Stealth[length=1.6mm]},
    data/.style={draw, rounded corners=2mm, fill=white, minimum height=7mm,
                 minimum width=16mm, align=center, inner sep=2pt},
    stage/.style={draw, fill=ln-box, minimum height=9mm, minimum width=22mm,
                  align=center, inner sep=2pt}]
  \node[data] (in) at (8,10) {\iflangja{3D\\プリミティブ}{3D\\primitives}};
  \node[stage, right=4mm of in] (geo) {\iflangja{ジオメトリ\\ステージ}{Geometry\\stage}};
  \node[data, right=4mm of geo] (mid) {\iflangja{2D\\プリミティブ}{2D\\primitives}};
  \node[stage, right=4mm of mid] (ras) {\iflangja{ラスタ化\\ステージ}{Rasterisation\\stage}};
  \node[data, right=4mm of ras, minimum width=12mm] (out) {\iflangja{ピクセル}{pixels}};
  \draw[->] (in) -- (geo);
  \draw[->] (geo) -- (mid);
  \draw[->] (mid) -- (ras);
  \draw[->] (ras) -- (out);
  \node[below=1.5mm of geo, text=ln-muted, align=center]
    {\iflangja{頂点シェーダ}{vertex shader}};
  \node[below=1.5mm of ras, text=ln-muted, align=center]
    {\iflangja{ピクセル（フラグメント）\\シェーダ}{pixel (fragment)\\shader}};
\end{tikzpicture}

% Width: body (120 mm); the drawing is about 116 mm wide.
\begin{tikzpicture}[x=1mm, y=1mm, font=\scriptsize\sffamily,
    >={Stealth[length=1.6mm]},
    data/.style={draw, rounded corners=2mm, fill=white, minimum height=7mm,
                 minimum width=16mm, align=center, inner sep=2pt},
    stage/.style={draw, fill=ln-box, minimum height=9mm, minimum width=22mm,
                  align=center, inner sep=2pt}]
  \node[data] (in) at (8,10) {\iflangja{3D\\プリミティブ}{3D\\primitives}};
  \node[stage, right=4mm of in] (geo) {\iflangja{ジオメトリ\\ステージ}{Geometry\\stage}};
  \node[data, right=4mm of geo] (mid) {\iflangja{2D\\プリミティブ}{2D\\primitives}};
  \node[stage, right=4mm of mid] (ras) {\iflangja{ラスタ化\\ステージ}{Rasterisation\\stage}};
  \node[data, right=4mm of ras, minimum width=12mm] (out) {\iflangja{ピクセル}{pixels}};
  \draw[->] (in) -- (geo);
  \draw[->] (geo) -- (mid);
  \draw[->] (mid) -- (ras);
  \draw[->] (ras) -- (out);
  \node[below=1.5mm of geo, text=ln-muted, align=center]
    {\iflangja{頂点シェーダ}{vertex shader}};
  \node[below=1.5mm of ras, text=ln-muted, align=center]
    {\iflangja{ピクセル（フラグメント）\\シェーダ}{pixel (fragment)\\shader}};
\end{tikzpicture}

% Width: body (120 mm); the drawing is about 116 mm wide.
\begin{tikzpicture}[x=1mm, y=1mm, font=\scriptsize\sffamily,
    >={Stealth[length=1.6mm]},
    data/.style={draw, rounded corners=2mm, fill=white, minimum height=7mm,
                 minimum width=16mm, align=center, inner sep=2pt},
    stage/.style={draw, fill=ln-box, minimum height=9mm, minimum width=22mm,
                  align=center, inner sep=2pt}]
  \node[data] (in) at (8,10) {\iflangja{3D\\プリミティブ}{3D\\primitives}};
  \node[stage, right=4mm of in] (geo) {\iflangja{ジオメトリ\\ステージ}{Geometry\\stage}};
  \node[data, right=4mm of geo] (mid) {\iflangja{2D\\プリミティブ}{2D\\primitives}};
  \node[stage, right=4mm of mid] (ras) {\iflangja{ラスタ化\\ステージ}{Rasterisation\\stage}};
  \node[data, right=4mm of ras, minimum width=12mm] (out) {\iflangja{ピクセル}{pixels}};
  \draw[->] (in) -- (geo);
  \draw[->] (geo) -- (mid);
  \draw[->] (mid) -- (ras);
  \draw[->] (ras) -- (out);
  \node[below=1.5mm of geo, text=ln-muted, align=center]
    {\iflangja{頂点シェーダ}{vertex shader}};
  \node[below=1.5mm of ras, text=ln-muted, align=center]
    {\iflangja{ピクセル（フラグメント）\\シェーダ}{pixel (fragment)\\shader}};
\end{tikzpicture}

  \caption{The traditional 3D graphics pipeline: geometry processing turns
    3D primitives into 2D primitives, rasterisation turns them into
    pixels.}
  \label{fig:l01-graphics-pipeline}
\end{figure}

The shaders were at first fixed-function units, then became
programmable, and finally the separate vertex and pixel units were
replaced by one pool of \term{unified shader}[統合型シェーダ]s, processors
that run either kind of shader program.  The unified shader is the
hallmark of NVIDIA's Tesla architecture, and with it the GPU became
usable for computation other than graphics.\margin{K12 §1.1 places the
inflection point at the G80 family (GeForce~8800, 2006), a programmable
unified graphics-and-compute processor, together with CUDA
\cite{kim2012gpgpu}.}  The term \term{GPGPU}[GPGPU] (general-purpose
computing on GPUs) was coined to mark that non-graphics use; now that
most GPU cycles in data centres go to non-graphics work the
distinction has lost its purpose, although the term survives in
titles.\sidenote{Before CUDA, GPGPU meant expressing a computation as a
graphics operation: data stored in textures, computation written as a
pixel shader in a shading language through DirectX or OpenGL, with
stream extensions such as Brook adding a thin layer on top (K12 §1.1
\cite{kim2012gpgpu}).  The unified shader and CUDA removed that
indirection.}

\subsection{Evolution of the programmable GPU}

Once programmable, the GPU acquired, generation by generation, the
features that general-purpose and then machine-learning workloads
needed:
\begin{itemize}
  \item cache hierarchies (L1 and L2 caches), where the graphics pipeline
    had needed only streaming access to textures;
  \item \term{double precision (FP64)}[倍精度]\index{FP64|see{double precision (FP64)}} arithmetic in addition to
    \term{single precision (FP32)}[単精度]\index{FP32|see{single precision (FP32)}}, for high-performance
    computing;\margin{FP64 at half the FP32 rate is a data-centre feature: the
    A100 delivers \num{9.7}~TFLOP/s FP64 against \num{19.5}~TFLOP/s
    FP32 \cite{nvidia2020ampere}.}
  \item atomic operations and fast integer arithmetic, for more diverse
    workloads;
  \item \term{high bandwidth memory (HBM)}[広帯域メモリ（HBM）]\index{HBM|see{high bandwidth memory (HBM)}}, DRAM dies
    stacked beside the GPU on an interposer and connected through a
    very wide interface, which raised memory bandwidth far beyond what
    a board-level memory bus provides;
  \item \term{half precision (FP16)}[半精度]\index{FP16|see{half precision (FP16)}}, a 16-bit floating-point
    format that doubles arithmetic throughput and halves memory traffic
    for machine-learning workloads that tolerate the reduced precision;
  \item the \term{tensor core}[テンソルコア], an execution unit that
    performs a small matrix multiply-add (for instance a
    $4 \times 4 \times 4$ product in one step) rather than a scalar
    operation per lane, for the dense linear algebra of neural
    networks;
  \item the \term{Transformer Engine}[Transformer Engine], Hopper's
    combination of FP8 tensor-core arithmetic with software and hardware
    that choose the number format and rescale layer by layer
    \cite{nvidia2022hopper}, for transformer models.\caution{The lecture says \enquote{transformer cores}.
    Hopper has no separate core for transformers: the Transformer
    Engine is the tensor cores' FP8 path plus the per-layer format
    management \cite{nvidia2022hopper}.}
\end{itemize}

\subsection{A modern GPU: NVIDIA H100}

The H100 (Hopper, 2022) illustrates the current state of the list.  It
supports FP64, FP32, FP16, INT8 and FP8 arithmetic; its tensor cores
and Transformer Engine serve machine learning; its register files are
larger;\caution{Lecture vs whitepaper: the register file per SM is
\SI{256}{\kibi\byte} on A100 and H100 alike; only the total grows with
the SM count, from \SI{27}{\mebi\byte} to \SI{33}{\mebi\byte} (HWP
Table~3 \cite{nvidia2022hopper}).} a \term{Tensor Memory Accelerator (TMA)}[テンソルメモリアクセラレータ]\index{TMA|see{Tensor Memory Accelerator (TMA)}}
moves multidimensional blocks of data between global and shared memory
asynchronously, relieving the threads of address generation; the number
of SMs and of floating-point units per SM has grown again; and an
NVLink Switch System connects many GPUs into one
machine.\margin{Lesson~12 treats FP8 and the tensor-core formats;
Lesson~9 the NVLink interconnect and multi-GPU systems.}
\Cref{tab:l01-generations} follows four data-centre generations through
the figures that this evolution changed.\caution{The H100 column holds
the preliminary figures of the V1.02 whitepaper
\cite{nvidia2022hopper}; shipped products differ in detail (the PCIe
H100 has 114 SMs).}

\begin{table}[htbp]
  \centering
  \caption{Four NVIDIA data-centre generations (A100 whitepaper
    \cite{nvidia2020ampere}; H100 whitepaper V1.02, SXM5, preliminary
    \cite{nvidia2022hopper}).}
  \label{tab:l01-generations}
  \setlength{\tabcolsep}{4pt}
  \begin{tabular}{>{\raggedright\arraybackslash}p{30mm}rrrr}
    \toprule
    & P100 & V100 & A100 & H100 \\
    \midrule
    Year & 2016 & 2017 & 2020 & 2022 \\
    SMs & 56 & 80 & 108 & 132 \\
    FP32 lanes per SM & 64 & 64 & 64 & 128 \\
    Tensor cores per SM & --- & 8 & 4 & 4 \\
    Register file per SM & \SI{256}{\kibi\byte} & \SI{256}{\kibi\byte} & \SI{256}{\kibi\byte} & \SI{256}{\kibi\byte} \\
    L2 cache & \SI{4}{\mebi\byte} & \SI{6}{\mebi\byte} & \SI{40}{\mebi\byte} & \SI{50}{\mebi\byte} \\
    \addlinespace
    HBM bandwidth & \SI{720}{\giga\byte\per\second} & \SI{900}{\giga\byte\per\second} & \SI{1555}{\giga\byte\per\second} & \SI{3000}{\giga\byte\per\second} \\
    Transistors & \num{15.3e9} & \num{21.1e9} & \num{54.2e9} & \num{80e9} \\
    Process & \SI{16}{\nano\metre} & \SI{12}{\nano\metre} & \SI{7}{\nano\metre} & 4N \\
    \bottomrule
  \end{tabular}
\end{table}

\begin{keypoint}[Lesson summary]
  A CPU minimises the latency of one thread with ILP, deep pipelines
  and large caches; a GPU maximises throughput with many small
  multithreaded cores, acts as an accelerator under a host CPU, and is
  programmed in the SPMD style on SIMT hardware.  Amdahl's law caps the
  speedup at $1/S$, so a parallel program must shrink its serial
  fraction; data decomposition with a reduction is the pattern that
  scales.  The GPU's core is the streaming multiprocessor, its lanes
  are stream processors, and the warp (32 threads on NVIDIA GPUs) is the
  unit of fetch and scheduling.  The GPU grew from a fixed-function
  graphics pipeline through programmable and unified shaders to a
  general processor with caches, FP64, HBM, FP16, tensor cores and a
  Transformer Engine.
\end{keypoint}

\end{document}
